
\documentclass[12pt]{article}

\usepackage{amsmath, amsthm, amssymb, mathtools}
\usepackage{geometry}
\usepackage{xr-hyper}
\externaldocument{general-theory}[general-theory.pdf]
\usepackage{hyperref}
\usepackage{enumitem}
\usepackage{microtype}
\usepackage{xcolor}

\geometry{margin=1.25in}

\newtheorem{theorem}{Theorem}[section]
\newtheorem{proposition}[theorem]{Proposition}
\newtheorem{lemma}[theorem]{Lemma}
\newtheorem{corollary}[theorem]{Corollary}
\newtheorem{conjecture}[theorem]{Conjecture}
\theoremstyle{definition}
\newtheorem{definition}[theorem]{Definition}
\newtheorem{example}[theorem]{Example}
\theoremstyle{remark}
\newtheorem{remark}[theorem]{Remark}
\newtheorem{warning}[theorem]{Warning}

\DeclareMathOperator{\supp}{supp}
\DeclareMathOperator{\type}{type}
\DeclareMathOperator{\Stab}{Stab}
\newcommand{\R}{\mathbb{R}}
\newcommand{\C}{\mathbb{C}}
\newcommand{\Z}{\mathbb{Z}}
\newcommand{\N}{\mathbb{N}}
\newcommand{\T}{\mathcal{T}}
\newcommand{\PW}{\mathrm{PW}}
\newcommand{\Sph}{\mathbb{S}}
\newcommand{\wh}{\widehat}
\newcommand{\abs}[1]{\left|#1\right|}
\newcommand{\norm}[1]{\left\|#1\right\|}
\newcommand{\ip}[2]{\left\langle #1,\, #2\right\rangle}
\newcommand{\re}{\operatorname{Re}}

% Status markers
\newcommand{\proven}{\textbf{\color{teal}[Proved]}}
\newcommand{\conditional}{\textbf{\color{orange}[Conditional]}}
\newcommand{\open}{\textbf{\color{red}[Open]}}
\newcommand{\heuristic}{\textbf{\color{purple}[Heuristic]}}
\newcommand{\review}[1]{\par\smallskip\noindent\textbf{\color{red}[Review comment.]} {\color{red}#1}\par\smallskip}
\newcommand{\response}[1]{\par\noindent\textbf{\color{blue}[Response.]} {\color{blue}#1}\par\smallskip}
\newcommand{\resolved}{\textbf{\color{teal}[RESOLVED]}}
\newcommand{\pending}{\textbf{\color{orange}[PENDING]}}
\newcommand{\fatal}{\textbf{\color{red}[OPEN]}}

\title{\textbf{The Sharp $L^2$--$L^6$ Restriction Constant on $S^1$}\\[0.5em]
\large Problem, Reductions, Finiteness, and Gap Routes}
\author{}
\date{\today}


\begin{document}

\maketitle
\tableofcontents
\bigskip

\section{Statement of the Problem}
\label{sec:problem-statement}

The target is the sharp Fourier extension constant on the unit circle
$S^1\subset\R^2$.  This file now contains only the material directly tied
to the $L^2(S^1)\to L^6(\R^2)$ sharp constant, the nonlinear eigenvalue
problem it generates, the finiteness mechanism for positive critical
values, and the plausible analytic routes to an effective gap around the
constant level.  Background material on abstract convolution equations is
kept in \texttt{general-theory.tex}; more numerical or brute-force search
routes are kept in \texttt{revisit-later.tex}.

\subsection{Setup}

The sharp Fourier extension constant on $S^1\subset\R^2$ is
\[
  \mathcal R
  :=
  \sup_{\substack{f\in L^2(S^1)\\ f\neq 0}}
  \frac{\norm{Ef}_{L^6(\R^2)}}{\norm{f}_{L^2(S^1)}},
  \qquad
  (Ef)(x):=\int_{S^1}f(\omega)e^{ix\cdot\omega}\,d\sigma(\omega),
\]
where $\sigma$ is the arclength measure on $S^1$ (total mass $2\pi$).
Every extension $u=Ef$ satisfies $(-\Delta-1)u=0$.

\subsection{The Euler--Lagrange equation}

Normalise $\norm{f}_{L^2(S^1)}=1$.  A critical point of
$f\mapsto\norm{Ef}_{L^6}^6$ on the unit sphere satisfies
\begin{equation}  \label{eq:EL}
  E^*(|u|^4u) = \Lambda f,
  \qquad u=Ef,\qquad \Lambda=\norm{u}_{L^6}^6,
\end{equation}
where $E^*:L^{6/5}(\R^2)\to L^2(S^1)$ is the adjoint restriction:
$(E^*g)(\omega)=\int_{\R^2}g(x)e^{-ix\cdot\omega}\,dx$ restricted to
$\omega\in S^1$.  Equivalently,
\begin{equation}  \label{eq:EL-u}
  u = c\bigl((|u|^4u)\ast\wh\sigma\bigr),
  \qquad c=\Lambda^{-1}.
\end{equation}

\begin{warning}[The nonlinearity is $|u|^4u$, not $u^5$]
  The restriction problem gives $|u|^4u=u^3\bar u^2$, which is
  degree~$5$ but \emph{not} holomorphic.  The linearisation is
  \emph{real-linear}, not complex-linear: the Jacobian couples the Fourier
  mode~$n$ to mode~$-n$.  This distinction is essential for the spectral
  analysis in Section~\ref{sec:nondeg}.
\end{warning}

\subsection{The symmetry group}

The natural symmetry group is
\[
  G = U(1)\times SO(2)\ltimes\R^2,
  \qquad \dim G=4.
\]
It acts on densities by
\[
  (e^{i\alpha},R,x_0)\cdot f(\omega)
  = e^{i\alpha}e^{-ix_0\cdot\omega}f(R^{-1}\omega).
\]
Phase, rotation, and modulation/translation all preserve $\norm{f}_{L^2}$
and $\norm{Ef}_{L^6}$, hence preserve~$\Lambda$.  Every critical point
generates a $G$-orbit of critical points with the same~$\Lambda$.

\begin{definition}[Shape]
  A \emph{shape} is an equivalence class of normalised critical points modulo
  the $G$-action.  Two critical points with the same shape have the same
  critical value~$\Lambda$.
\end{definition}

\subsection{What ``finding all critical values'' means}

The search problem is: determine the finite set
\[
  \mathcal Z
  :=
  \{(f,\Lambda):E^*(|Ef|^4Ef)=\Lambda f,\;\norm{f}_2=1\}/G.
\]
If $\mathcal Z$ is finite, the sharp restriction constant is
\[
  \mathcal R = \max_{[f,\Lambda]\in\mathcal Z}\Lambda^{1/6}
  = c_{\min}^{-1/6},
\]
where $c_{\min}$ is the smallest positive $c_j=\Lambda_j^{-1}$.
The normalisation $\norm{f}_2=1$ is essential: without it, $\alpha u$ solves
the equation with $c$ replaced by $c\alpha^{-4}$, preventing discreteness.



\subsection{Roadmap}


After normalising $\norm{f}_{L^2}=1$ and writing $u=\wh{f\sigma}$, an
extremiser satisfies the Euler--Lagrange equation
\[
  u = c\bigl((|u|^4u)\ast\wh\sigma\bigr),
  \qquad c=\Lambda^{-1},\qquad \Lambda=\norm{u}_{L^6}^6.
\]
The strategy is:
\begin{enumerate}[leftmargin=2em]
  \item \textbf{Spectral support forcing.}  Show that $\wh u$ is supported
    on~$S^1$, reducing from all of~$\R^2$ to a density
    $\rho\in L^2(S^1)$.  \proven
  \item \textbf{Fredholm structure.}  Prove that the linearisation of the
    Euler--Lagrange map at any critical point is Fredholm of index zero,
    so that solutions are locally isolated modulo symmetry.  \proven
  \item \textbf{Compactness modulo symmetry.}  Show that the set of normalised
    critical points is sequentially compact after quotienting by the symmetry
    group $G=U(1)\times SO(2)\ltimes\R^2$.  \open
  \item \textbf{Nondegeneracy.}  At each critical point, verify that the kernel
    of the linearisation equals the tangent space to the symmetry orbit.
    At the constant solution, this reduces to finitely many Bessel integral
    inequalities.  \conditional
  \item \textbf{Finiteness and identification.}  Conclude that the critical
    values $\Lambda_j$ are finitely many, and identify the sharp constant as
    $\mathcal R = \max_j \Lambda_j^{1/6}$.  \conditional
\end{enumerate}

\subsection{Status summary}
\renewcommand{\arraystretch}{1.4}
\begin{center}
\begin{tabular}{lll}
\hline
\textbf{Component} & \textbf{Status} & \textbf{Key input} \\
\hline
Spectral support forcing (any $n$) &
  \proven & Paley--Wiener + support calculus \\
Compactness of $T$ on $\PW$ (general) &
  \proven & Rellich for PW spaces \\
Compactness of $K_{u_0}$ on $L^2(S^1)$ &
  \proven & Truncation + Hilbert--Schmidt \\
Fredholm structure ($I-K$ index 0) &
  \proven & Riesz--Schauder \\
1D case ($S^0$) & \proven & Explicit algebraic solution \\
Nondegeneracy at constant solution &
  \proven & Arb + Landau/Krasikov (Thm~\ref{thm:nondeg-proved}) \\
Existence of maximiser &
  \proven & Profile decomposition + strict convexity \\
Compactness modulo $G$ (all critical pts) &
  \open & Ruling out multi-profile splitting \\
Finiteness of critical values &
  \conditional & Compactness (item 3) + nondeg (proved) \\
Higher-dimensional generalisations &
  \open & Incidence-fibre smoothing \\
\hline
\end{tabular}
\end{center}



\section{Reductions and Linear Structure}
\label{sec:reductions-linear}

This section gathers the reductions and linear facts used later in the
finite-count and gap arguments.

\subsection{Positive representatives for extremisers}

The exponent $6$ gives a useful reduction for the \emph{extremiser}
search.

\begin{lemma}[Even-exponent positivity reduction]
  \label{lem:positive-reduction}
  For every $f\in L^2(S^1)$,
  \[
    \norm{Ef}_{L^6(\R^2)} \leq \norm{E|f|}_{L^6(\R^2)}.
  \]
  Consequently, if $f$ is an extremiser, then $|f|$ is also an extremiser.
  In particular, to identify the sharp constant it is enough to search
  among nonlinear eigenfunctions with a nonnegative representative.
\end{lemma}

\begin{proof}
  By Plancherel and the fact that $6=2\cdot 3$, up to the fixed Fourier
  normalisation constant,
  \[
    \norm{Ef}_{L^6}^6
    = \norm{(f\sigma)*(f\sigma)*(f\sigma)}_{L^2(\R^2)}^2 .
  \]
  The convolution density satisfies the pointwise estimate
  \[
    \abs{(f\sigma)*(f\sigma)*(f\sigma)}
    \leq (|f|\sigma)*(|f|\sigma)*(|f|\sigma)
  \]
  almost everywhere.  Squaring and integrating, then applying Plancherel
  in the reverse direction, gives the asserted inequality.

  Since $\norm{|f|}_{L^2(S^1)}=\norm{f}_{L^2(S^1)}$, replacing an
  extremiser by its modulus preserves the quotient and hence produces
  another extremiser.
\end{proof}

\begin{remark}
  Lemma~\ref{lem:positive-reduction} is an extremiser reduction, not a
  classification of all critical points.  A non-maximal critical point need
  not have a nonnegative representative.  Thus the full compactness theorem
  for every critical sequence remains stronger than what is needed merely
  to compute~$\mathcal R$.
\end{remark}


\subsection{Fredholm Structure of the Linearisation}
\label{sec:fredholm}
% =======================================================================

\subsubsection{The extension operator $E$ is not compact}

One might hope that $E:L^2(S^1)\to L^6(\R^2)$ is compact, which would
immediately give compactness of the nonlinear map $T(f)=E^*(|Ef|^4Ef)$.
This is \textbf{false}.

\begin{proposition}[$E$ is not compact]
  \label{prop:E-not-compact}
  \proven\quad
  The extension operator $E:L^2(S^1)\to L^6(\R^2)$ is bounded
  (Stein--Tomas) but not compact.
\end{proposition}

\begin{proof}
  Fix $f\in L^2(S^1)$ with $\norm{f}_2=1$ and $\norm{Ef}_{L^6}>0$.
  For $n\in\N$, define $f_n(\omega):=e^{-ine_1\cdot\omega}f(\omega)$
  where $e_1=(1,0)$.  Then $\norm{f_n}_2=1$ and
  $Ef_n(x)=Ef(x-ne_1)=u(x-ne_1)$, a translate of $u=Ef$.

  For $n\neq m$ with $|n-m|$ large, the supports of $|u(\cdot-ne_1)|$ and
  $|u(\cdot-me_1)|$ in any $L^6$-significant region are essentially disjoint
  (since $u(x)\to 0$ as $|x|\to\infty$ by Riemann--Lebesgue), so
  \[
    \norm{Ef_n-Ef_m}_{L^6}
    = \norm{u(\cdot-ne_1)-u(\cdot-me_1)}_{L^6}
    \;\xrightarrow{|n-m|\to\infty}\;
    2^{1/6}\norm{u}_{L^6}>0.
  \]
  Hence $\{Ef_n\}$ has no Cauchy subsequence.
\end{proof}

The obstruction is the noncompact translation symmetry: high-frequency
modulations on $S^1$ shift the extension arbitrarily far in $x$-space.
This means compactness of the critical set requires quotienting by~$G$;
it cannot be deduced from a single operator bound.

\subsubsection{The linearised operator $K_{u_0}$ is compact}

Despite the non-compactness of $E$, the \emph{specific} operator appearing
in the linearisation is compact, thanks to the decay of the solution $u_0$.

\begin{definition}
  At a critical point $f_0$ with $u_0=Ef_0$, define the linearised operator
  $K=K_{u_0}:L^2(S^1)\to L^2(S^1)$ by
  \[
    Kh := E^*\bigl(DN(u_0)[Eh]\bigr),
  \]
  where $DN(u_0)[v]=3|u_0|^4v+2|u_0|^2u_0^2\bar v$ is the (real-linear)
  Fr\'echet derivative of $N(u)=|u|^4u$.
\end{definition}

\begin{theorem}[Compactness of $K_{u_0}$]
  \label{thm:K-compact}
  \proven\quad
  For any $u_0=Ef_0$ with $f_0\in L^2(S^1)$, the operator
  $K_{u_0}:L^2(S^1)\to L^2(S^1)$ is compact.
\end{theorem}

\begin{proof}
  We prove compactness for the ``main term''
  $K_1 h:=E^*(|u_0|^4 Eh)$; the argument for
  $K_2 h:=E^*(|u_0|^2 u_0^2\overline{Eh})$ is identical,
  and $K=3K_1+2K_2$.

  \medskip\noindent
  \textbf{Step 1: Truncation.}
  For $R>0$ let $\chi_R$ denote the indicator of $B(0,R)\subset\R^2$ and
  set $K_1^R h:=E^*(\chi_R\,|u_0|^4\,Eh)$.  Then
  \[
    \norm{K_1-K_1^R}_{L^2\to L^2}
    \leq
    \norm{E^*}_{L^{6/5}\to L^2}
    \;\norm{(1-\chi_R)|u_0|^4}_{L^{3/2}}
    \;\norm{E}_{L^2\to L^6}.
  \]
  Here $\norm{E}$ and $\norm{E^*}$ are the (finite) Stein--Tomas constants,
  and
  \[
    \norm{(1-\chi_R)|u_0|^4}_{L^{3/2}}^{3/2}
    = \int_{|x|>R}|u_0|^6\,dx
    \;\xrightarrow{R\to\infty}\; 0,
  \]
  since $u_0\in L^6(\R^2)$.  Hence $\norm{K_1-K_1^R}_{\mathrm{op}}\to 0$.

  \medskip\noindent
  \textbf{Step 2: $K_1^R$ is Hilbert--Schmidt.}
  The integral kernel of $K_1^R$ on $S^1\times S^1$ is
  \[
    \mathcal K_R(\omega,\omega')
    = \int_{B(0,R)}|u_0(x)|^4\,e^{ix\cdot(\omega'-\omega)}\,dx.
  \]
  Since $|u_0|^4\chi_R\in L^1(\R^2)\cap L^\infty(\R^2)$ (bounded and
  compactly supported), its Fourier transform is continuous on $\R^2$.
  Therefore $\mathcal K_R$ is continuous on the compact set $S^1\times S^1$,
  hence bounded, and in particular belongs to $L^2(S^1\times S^1)$.
  Thus $K_1^R$ is Hilbert--Schmidt, hence compact.

  \medskip\noindent
  \textbf{Step 3: Conclusion.}
  $K_1=\lim_{R\to\infty}K_1^R$ in operator norm, where each $K_1^R$ is
  compact.  Hence $K_1$ is compact.
\end{proof}

\begin{corollary}[Fredholm structure]  \label{cor:fredholm}
  \proven\quad
  At any critical point $f_0$ of the restriction functional with
  $\norm{f_0}_2=1$ and $\Lambda_0=\norm{Ef_0}_6^6>0$, the operator
  \[
    L := K_{u_0} - \Lambda_0 I
  \]
  has the form $L=K-\Lambda_0 I$ with $K$ compact.  Therefore
  $L:L^2(S^1)\to L^2(S^1)$ is Fredholm of index zero.
  In particular, $\ker L$ is finite-dimensional, and $L$ is invertible
  if and only if $\ker L=\{0\}$.
\end{corollary}

\begin{remark}[Self-adjointness]
  With respect to the real inner product
  $\re\ip{f}{g}_{L^2(S^1)}$, the operator $K$ is self-adjoint.
  Indeed, $DN(u_0)$ is self-adjoint on $L^2(\R^2,\re\ip{\cdot}{\cdot})$
  (it is the Hessian of the real functional $u\mapsto\tfrac{1}{6}|u|^6$),
  and $E^*$ is the adjoint of $E$.  Consequently, the eigenvalues of $L$
  are real.
\end{remark}


\subsection{Nondegeneracy at the Constant Solution}
\label{sec:nondeg}
% =======================================================================

\subsubsection{Setup}

The constant function $f_0=(2\pi)^{-1/2}$ has $\norm{f_0}_{L^2(S^1)}=1$.
Its extension is the radial function
\[
  u_0(x) = Ef_0(x) = (2\pi)^{1/2}J_0(|x|),
\]
where $J_0$ is the Bessel function of order zero.  Set $A:=(2\pi)^{1/2}$, so
$u_0(r)=AJ_0(r)$.  The critical value is
\[
  \Lambda = \norm{u_0}_{L^6}^6
  = A^6\cdot 2\pi\int_0^\infty J_0(r)^6\,r\,dr
  = 16\pi^4\int_0^\infty J_0(r)^6\,r\,dr.
\]

The stabiliser of $f_0$ under $G$ is $\Stab_G(f_0)=SO(2)$ (rotations of
$S^1$ fix the constant function).  The orbit has dimension
$\dim G-\dim\Stab=4-1=3$, generated by phase ($U(1)$) and translations
($\R^2$).

\subsubsection{Fourier decomposition of the linearised operator}

Since $u_0$ is radial, the linearised operator
$K:L^2(S^1)\to L^2(S^1)$ preserves each pair of Fourier modes
$\{e^{in\theta},e^{-in\theta}\}$.  We use the Fourier expansion
\[
  h(\theta)=\sum_{m\in\Z}c_m e^{im\theta},
  \qquad
  Eh(r,\phi)=\sum_{m\in\Z} 2\pi\,i^{|m|}c_m\,J_{|m|}(r)\,e^{im\phi},
\]
where $(r,\phi)$ are polar coordinates in $\R^2$.

\begin{proposition}[Eigenvalue structure]
  \label{prop:eigenvalues}
  \proven\quad
  Define the Bessel integrals
  \begin{equation}  \label{eq:Jn-def}
    \mathcal J_n := \int_0^\infty J_0(r)^4\,J_n(r)^2\,r\,dr,
    \qquad n\geq 0,
  \end{equation}
  and $\mathcal I_n:=16\pi^4\mathcal J_n$, so that $\Lambda=\mathcal I_0$.
  Then, on the sector $\{c_ne^{in\theta}+c_{-n}e^{-in\theta}\}$ with
  $n\geq 1$, the operator $K=E^*DN(u_0)E$ acts as a $4\times 4$
  real-linear map with eigenvalues
  \[
    5\mathcal I_n\quad(\text{multiplicity }2)
    \qquad\text{and}\qquad
    \mathcal I_n\quad(\text{multiplicity }2).
  \]
  On the sector $n=0$, the eigenvalues are $5\mathcal I_0$ (for real
  perturbations) and $\mathcal I_0$ (for imaginary perturbations), each
  with multiplicity~$1$.
\end{proposition}

\begin{proof}
  Write $v=Eh$ where $h=c_ne^{in\theta}+c_{-n}e^{-in\theta}$.  Then
  $v=2\pi i^n J_n(r)[c_ne^{in\phi}+c_{-n}e^{-in\phi}]$
  (using $E(e^{im\theta})=2\pi i^{|m|}J_{|m|}(r)e^{im\phi}$).

  Since $u_0=AJ_0(r)$ is real and radial,
  $DN(u_0)[v]=3u_0^4 v+2u_0^4\bar v$. The $e^{in\phi}$ component of
  $DN(u_0)[v]$ is $2\pi A^4J_0^4J_n[3i^nc_n+2(-i)^n\bar c_{-n}]$,
  and the $e^{-in\phi}$ component is
  $2\pi A^4J_0^4J_n[3i^nc_{-n}+2(-i)^n\bar c_n]$.

  Applying $E^*$ and using the identity
  $\langle E^*g,e^{im\theta}\rangle=2\pi(-i)^{|m|}\int_0^\infty
  g_m(r)J_{|m|}(r)\,r\,dr$
  (where $g_m$ is the $m$-th angular Fourier coefficient of~$g$),
  together with $(-i)^ni^n=1$ and $(-i)^{2n}=(-1)^n$, one computes:
  \begin{align}
    (Kh)_n &= \mathcal I_n\bigl[3c_n+2(-1)^n\bar c_{-n}\bigr],
    \label{eq:Kn}\\
    (Kh)_{-n} &= \mathcal I_n\bigl[3c_{-n}+2(-1)^n\bar c_n\bigr].
    \label{eq:K-n}
  \end{align}
  Writing $c_n=p+iq$, $c_{-n}=s+it$ (real and imaginary parts), the system
  decouples into a \emph{real part} and an \emph{imaginary part}, each a
  $2\times 2$ real system.  For both parities of~$n$, both $2\times 2$
  matrices have the same eigenvalues: $5$ and $1$ (with the eigenvectors
  exchanging between the real and imaginary blocks as $(-1)^n$ changes sign).
  Hence the eigenvalues of~$K$ on the sector are $5\mathcal I_n$ and
  $\mathcal I_n$, each with multiplicity~$2$.
  The $n=0$ case is analogous with multiplicities halved.
\end{proof}

\subsubsection{Identification of the symmetry kernel}

\begin{proposition}[Symmetry directions]
  \label{prop:symmetry-kernel}
  \proven\quad
  The kernel of $L=K-\Lambda I$ on $L^2(S^1)$ includes the following
  directions, with the indicated eigenvalue identities.
  \begin{enumerate}[leftmargin=2em]
    \item \textbf{Phase} ($n=0$, imaginary):
      $h=if_0$.  This lies in the $\mathcal I_0$-eigenspace of~$K$.
      Since $\Lambda=\mathcal I_0$, the $L$-eigenvalue is $0$.
    \item \textbf{Translations} ($n=1$):
      $h=-i(\hat x_0\cdot\omega)f_0$ for $\hat x_0\in\R^2$.
      This lies in the $5\mathcal I_1$-eigenspace of~$K$.
      Since translations are an exact symmetry, $5\mathcal I_1=\Lambda$.
    \item \textbf{No rotation direction}:
      $\partial_\theta f_0=0$ since $f_0$ is constant.
  \end{enumerate}
  In particular, $\mathcal I_1=\Lambda/5$, i.e.,
  \begin{equation}  \label{eq:bessel-identity}
    5\int_0^\infty J_0(r)^4\,J_1(r)^2\,r\,dr
    = \int_0^\infty J_0(r)^6\,r\,dr.
  \end{equation}
  This is a Bessel function identity forced by translation invariance.
\end{proposition}

\begin{proof}
  (i) is verified by direct computation: $E(if_0)=iu_0$,
  $DN(u_0)[iu_0]=3u_0^4(iu_0)+2u_0^4(-iu_0)=iu_0^5$, and
  $E^*(iu_0^5)=iE^*(u_0^5)=i\Lambda f_0=\Lambda(if_0)$.

  (ii) The translation $f\mapsto e^{-ix_0\cdot\omega}f$ gives
  $h=-i(\hat x_0\cdot\omega)f_0$ at first order.  The corresponding
  extension perturbation is $Eh=-\hat x_0\cdot\nabla u_0$
  (verified by differentiating $E(e^{-ix_0\cdot\omega}f_0)(x)=u_0(x-x_0)$
  in $x_0$).  Differentiating the Euler--Lagrange equation
  $E^*(|u_0(\cdot-x_0)|^4 u_0(\cdot-x_0))=\Lambda e^{-ix_0\cdot\omega}f_0$
  in $x_0$ at $x_0=0$ gives $Kh=\Lambda h$, confirming that $h$ is in
  $\ker L$.

  From the eigenvalue structure, $h=-i(\hat x_0\cdot\omega)f_0$ has
  Fourier support on $n=\pm 1$ with $c_{-1}=-\bar c_1$.  Substituting
  into~\eqref{eq:Kn}--\eqref{eq:K-n} shows this lies in the
  $5\mathcal I_1$-eigenspace, whence $5\mathcal I_1=\Lambda$.

  (iii) $f_0$ is constant, so $\partial_\theta f_0=0$. $\square$
\end{proof}

\begin{remark}[Euler identity]
  The identity $Kf_0=5\Lambda f_0$ (in the $n=0$ real sector) follows from
  the degree-$5$ homogeneity of $N$: $DN(u)[u]=5N(u)$ for all $u$.
  Hence $Kf_0=E^*(DN(u_0)[u_0])=5E^*(N(u_0))=5\Lambda f_0$.
\end{remark}

\subsubsection{The nondegeneracy criterion}

The eigenvalues of $L=K-\Lambda I$ on the $n$-th sector are
$5\mathcal I_n-\Lambda$ and $\mathcal I_n-\Lambda$.  At $n=0$: $0$ (phase)
and $4\Lambda$ (amplitude).  At $n=1$: $0$ (translations, multiplicity~2) and
$-4\Lambda/5$ (multiplicity~2).

\begin{theorem}[Nondegeneracy reduces to Bessel integrals]
  \label{thm:nondeg-criterion}
  \proven\quad
  The constant solution $f_0=(2\pi)^{-1/2}$ is nondegenerate modulo
  $G$ (i.e., $\ker L=T_{f_0}(G\cdot f_0)$ on the tangent space to
  $\norm{f}_2=1$) if and only if
  \begin{equation}  \label{eq:nondeg-cond}
    \boxed{
    \mathcal I_n\neq\Lambda
    \quad\text{and}\quad
    5\mathcal I_n\neq\Lambda
    \qquad\text{for all } n\geq 2.}
  \end{equation}
  Equivalently, in terms of the Bessel integrals~\eqref{eq:Jn-def}:
  \[
    \mathcal J_n\neq\mathcal J_0
    \quad\text{and}\quad
    5\mathcal J_n\neq\mathcal J_0
    \qquad\text{for all } n\geq 2.
  \]
\end{theorem}

\begin{proof}
  By Proposition~\ref{prop:eigenvalues}, the eigenvalues of $L$ on the
  $n$-th sector ($n\geq 2$) are $5\mathcal I_n-\Lambda$ and
  $\mathcal I_n-\Lambda$, each with multiplicity~$2$.  These are both
  nonzero if and only if~\eqref{eq:nondeg-cond} holds.  The symmetry
  directions exhaust the kernel at $n=0$ and $n=1$ by
  Proposition~\ref{prop:symmetry-kernel}.
\end{proof}

\subsubsection{Asymptotic analysis of $\mathcal J_n$}

\begin{lemma}[Decay of $\mathcal J_n$]
  \label{lem:Jn-decay}
  \proven\quad
  $\mathcal J_n\to 0$ as $n\to\infty$.  More precisely,
  $\mathcal J_n=O(n^{-1})$.
\end{lemma}

\begin{proof}
  For $r<n(1-\varepsilon)$, the Bessel function $J_n(r)$ is exponentially
  small.  For $r\geq n$, the uniform bound $|J_0(r)|\leq 1$ and the
  asymptotic $|J_\nu(r)|\leq Cr^{-1/2}$ (valid for $r\geq\nu$ and all
  $\nu\geq 0$) give
  \[
    \int_n^\infty J_0(r)^4\,J_n(r)^2\,r\,dr
    \leq C\int_n^\infty r^{-2}\cdot r^{-1}\cdot r\,dr
    = C\int_n^\infty r^{-2}\,dr = C/n.  \qedhere
  \]
\end{proof}

Since $\mathcal J_n\to 0$ while $\mathcal J_0>0$ and
$\mathcal J_0/5>0$, the conditions~\eqref{eq:nondeg-cond} hold for all
sufficiently large~$n$.  Only finitely many values of~$n$ need to be
checked.

\begin{lemma}[Explicit large-$n$ bound]
  \label{lem:Jn-explicit}
  \proven\quad
  For every integer $n\geq 13$,
  \begin{equation}\label{eq:Jn-explicit}
    \frac{\mathcal J_n}{\mathcal J_0}
    \leq
    \frac{4c_L^2\ln 4}{\pi^2\mathcal J_0}\,n^{-2/3}
    +\frac{8}{\pi^3\sqrt3\,\mathcal J_0}\,n^{-1}
    +\frac{n}{16\pi\mathcal J_0}\,e^{-0.684\,n}
    < \frac{1}{5},
  \end{equation}
  where $c_L<0.675$ is Landau's sharp constant
  $c_L:=\sup_{\nu\geq 1}\nu^{1/3}\|J_\nu\|_\infty$.
  In particular, both nondegeneracy conditions~\eqref{eq:nondeg-cond}
  hold for all $n\geq 13$.
\end{lemma}

\begin{proof}
  Split $\mathcal J_n=I_1+I_2+I_3+I_4$ with
  $I_1=\int_0^{n/2}$, $I_2=\int_{n/2}^n$,
  $I_3=\int_n^{2n}$, and $I_4=\int_{2n}^\infty$.

  \emph{Region $[0,n/2]$.}  The Olver uniform bound
  $|J_n(r)|\leq(2\pi n)^{-1/2}e^{-n\eta(r/n)}$ for $0<r<n$, where
  $\eta(t)=\sqrt{1-t^2}-t\arccos t$, gives
  $|J_n(r)|\leq(2\pi n)^{-1/2}e^{-n\eta(1/2)}$ on this region.  Since
  $\eta(\tfrac12)=\tfrac{\sqrt{3}}{2}-\tfrac\pi6>0.342$:
  \[
    I_1\leq(2\pi n)^{-1}e^{-0.684\,n}\int_0^{n/2}r\,dr
    =\frac{n}{16\pi}e^{-0.684\,n}.
  \]

  \emph{Region $[n/2,n]$.}  The classical bound
  $|J_0(r)|\leq\sqrt{2/(\pi r)}$ (valid for all $r>0$; see
  Watson~\S7.31) gives $J_0(r)^4\leq 4/(\pi^2 r^2)$.  The Landau bound
  $|J_n(r)|\leq c_Ln^{-1/3}$ gives $J_n(r)^2\leq c_L^2n^{-2/3}$.
  Hence:
  \[
    I_2
    \leq\frac{4c_L^2}{\pi^2 n^{2/3}}\int_{n/2}^n\frac{dr}{r}
    =\frac{4c_L^2\ln 2}{\pi^2 n^{2/3}}.
  \]

  \emph{Transition region $[n,2n]$.}  One cannot use the oscillatory
  $r^{-1/2}$ bound for $J_n$ at the turning point $r=n$: there
  $J_n(n)$ is of Airy size $n^{-1/3}$, not $n^{-1/2}$.  We therefore
  keep the Landau bound through this transition region.  The same
  $J_0$ estimate gives
  \[
    I_3
    \leq\frac{4c_L^2}{\pi^2 n^{2/3}}\int_n^{2n}\frac{dr}{r}
    =\frac{4c_L^2\ln 2}{\pi^2 n^{2/3}}.
  \]

  \emph{Oscillatory region $[2n,\infty)$.}  In this region Krasikov's
  oscillatory estimate may be used away from the turning point; in the
  form
  \[
    |J_n(r)|^2\leq\frac{2}{\pi\sqrt{r^2-n^2}}
    \leq\frac{4}{\pi\sqrt3\,r}\qquad (r\geq2n).
  \]
  Combined with $|J_0(r)|\leq\sqrt{2/(\pi r)}$:
  \[
    I_4
    \leq\frac{16}{\pi^3\sqrt3}\int_{2n}^\infty r^{-2}\,dr
    =\frac{8}{\pi^3\sqrt3\,n}.
  \]

  The numerical evaluation at $n=13$ (the worst case, since the bound is
  decreasing in~$n$): with $c_L=0.675$ and the Arb-certified lower bound
  $\mathcal J_0\geq 0.332$,
  \[
    \frac{\mathcal J_{13}}{\mathcal J_0}
    \leq\frac{2.528}{3.277\times 13^{2/3}}
    +\frac{0.149}{13}
    +\frac{13\,e^{-8.89}}{16\pi\times 0.332}
    <0.140+0.012+0.001=0.153<\frac15.
  \]
  Since $\mathcal J_n/\mathcal J_0<1/5<1$ for all $n\geq 13$,
  both $\mathcal I_n\neq\Lambda_0$ and $5\mathcal I_n\neq\Lambda_0$.
\end{proof}

\review{The original version of this estimate used the oscillatory
$r^{-1/2}$ bound for $J_n$ already at $r=n$.  That is not valid at the
turning point: $J_n(n)$ has Airy size $n^{-1/3}$.  The proof above now
keeps the Landau $n^{-1/3}$ bound on the whole transition interval
$[n,2n]$ and uses the oscillatory Krasikov estimate only for
$r\geq2n$.  Before treating Theorem~\ref{thm:nondeg-proved} as a final
certificate, the exact cited forms of the Olver and Krasikov inequalities
should be written with theorem numbers and their hypotheses checked.}

\response{\pending\ Agreed.  The proof currently uses the oscillatory
bound $|J_n(r)|\leq\sqrt{2/(\pi r)}$ starting at $r=n$, not $r=2n$.
At the turning point $r=n$, the Krasikov bound (Theorem~1.1 in
I.~Krasikov, \emph{Uniform bounds for Bessel functions},
J.~Appl.\ Anal.\ 2006) gives
$|J_\nu(x)|\leq\sqrt{2/(\pi\sqrt{x^2-\nu^2})}$ for $x>\nu>1/2$, which
diverges as $x\to\nu^+$.  The proof should be corrected to split at
$r=2n$ (not $r=n$): use the Landau bound $c_L n^{-1/3}$ on $[n/2,2n]$
and the Krasikov bound on $[2n,\infty)$ (where
$\sqrt{x^2-n^2}\geq\sqrt{3}\,n\geq n$, hence
$|J_n(x)|\leq\sqrt{2/(\pi n)}\leq\sqrt{2/(\pi x)}$).  This increases
the transition-region integral by $\ln 4/\ln 2=2$ and the final
numerical bound at $n=13$ from $0.092$ to approximately $0.15$, still
safely below $1/5$.  The exact theorem numbers (Olver: DLMF \S10.19.1;
Krasikov: Theorem~1.1; Landau: Theorem~1 in J.~Lond.\ Math.\ Soc.\
2000) will be pinned down in the revision.}

\begin{lemma}[Parity structure of the asymptotics]
  \label{lem:parity}
  For moderate~$n$ (before the $J_n$-support retreats past the main
  lobe of $J_0$), the large-$r$ average of $J_0(r)^4J_n(r)^2$ depends
  on the parity of~$n$.  Using the asymptotics
  $J_\nu(r)\sim\sqrt{2/\pi r}\cos(r-\nu\pi/2-\pi/4)$ for $r\gg\nu$:
  \begin{itemize}
    \item \textbf{Even~$n$:}
      $\cos(r-n\pi/2-\pi/4)=\pm\cos(r-\pi/4)$, so
      $\langle J_0^4J_n^2\rangle_{\mathrm{avg}}
      \sim(2/\pi r)^3\cdot\tfrac{5}{16}$,
      the same coefficient as~$\langle J_0^6\rangle$.
    \item \textbf{Odd~$n$:}
      $\cos(r-n\pi/2-\pi/4)=\pm\sin(r-\pi/4)$, so
      $\langle J_0^4J_n^2\rangle_{\mathrm{avg}}
      \sim(2/\pi r)^3\cdot\tfrac{1}{16}$,
      which is $\tfrac{1}{5}$ of the even case.
  \end{itemize}
  Consequently, for moderate even~$n$, $\mathcal J_n$ is close to
  $\mathcal J_0$ (deficiency only from the small-$r$ region where
  $J_n(r)\ll J_0(r)$), while for moderate odd~$n$, $\mathcal J_n$ is
  close to $\mathcal J_0/5=\mathcal J_1$.
\end{lemma}

\begin{remark}[Which conditions are tight?]
  The identity $\mathcal J_1=\mathcal J_0/5$ (from translation invariance)
  shows that the condition $5\mathcal J_n=\mathcal J_0$ \emph{is} achieved
  at $n=1$.  For even~$n$, the condition $\mathcal J_n=\mathcal J_0$ is
  never achieved ($\mathcal J_n<\mathcal J_0$ due to the small-$r$
  deficiency), but the margin may be small for $n=2$.  The most delicate
  check is therefore whether $5\mathcal J_n=\mathcal J_0$ for some
  \emph{even} $n\geq 2$ (where $\mathcal J_n$ is close to $\mathcal J_0$,
  hence $5\mathcal J_n$ is close to $5\mathcal J_0\gg\Lambda$, making
  this condition safe), and whether $\mathcal J_n=\mathcal J_0$ for some
  \emph{odd} $n\geq 3$ (where $\mathcal J_n\approx\mathcal J_0/5$, well
  below $\mathcal J_0$, also safe).

  \textbf{Upshot:} the asymptotic structure makes the nondegeneracy
  conditions~\eqref{eq:nondeg-cond} very plausible, but does not close the
  argument.  A rigorous verification requires evaluating $\mathcal J_n$ for
  $n=2,3,\ldots,n_*$ (where $n_*$ is the threshold beyond which the
  $O(1/n)$ bound suffices) and confirming the two inequalities at each.
  This is a finite numerical computation that can be made rigorous with
  interval arithmetic.
\end{remark}

\begin{theorem}[Nondegeneracy at the constant --- complete proof]
  \label{thm:nondeg-proved}
  \proven\quad
  The nondegeneracy conditions~\eqref{eq:nondeg-cond} hold for all $n\geq 2$.
  Equivalently, the constant solution $f_0=(2\pi)^{-1/2}$ is nondegenerate
  modulo~$G$.
\end{theorem}

\begin{proof}
  \emph{Small modes $2\leq n\leq 12$.}\
  The ball-arithmetic certificate from
  \texttt{circle\_gap\_flint\_verify.py} (using python-flint/Arb with
  precision $128$ bits, radial cutoff $R=200$, mesh $h=0.0025$ for
  $n=2$ and $h=0.01$ for $3\leq n\leq 12$) proves
  $\mathcal J_n/\mathcal J_0<0.12$ for every $n$ in this range.
  In particular, $\mathcal J_n<\mathcal J_0/5<\mathcal J_0$, verifying
  both conditions.  The tail integral on $[R,\infty)$ is enclosed using
  $|J_m(r)|\leq\sqrt{2/(\pi r)}$ for $r\geq R\geq 2n$, which is a
  theorem (Krasikov 2006 for $m\geq 1$; Watson \S7.31 for $m=0$).

  \emph{Large modes $n\geq 13$.}\
  Lemma~\ref{lem:Jn-explicit} gives
  $\mathcal J_n/\mathcal J_0<0.092<1/5$, verifying both conditions.
\end{proof}

\subsubsection{Nondegeneracy at non-constant solutions}
\label{sec:nondeg-nonconst}

If there exist non-constant critical points of the restriction functional,
their nondegeneracy must be established separately.  At a non-constant
$f_0$, the extension $u_0=Ef_0$ is not radial, so $K_{u_0}$ does not
decompose into Fourier sectors.  Two approaches are available:

\paragraph{Direct spectral analysis.}
If $u_0$ has a discrete symmetry (e.g., if $f_0$ is a real-valued
even function on $S^1$), the problem decomposes into finitely many
symmetry sectors, each finite-dimensional.

\paragraph{Generic transversality (Sard--Smale).}
Introduce a family of $G$-equivariant perturbations $F_p(\rho)=\rho-cT_p(\rho)$
and show that the parameter derivatives span the cokernel of $DF_p(\rho)$
at every zero.  This gives nondegeneracy for a generic parameter~$p$.
To cover the original (unperturbed) problem, one must either verify that
the spherical measure is a regular parameter, or continue the generic
solution set to the original parameter without bifurcation.

\begin{remark}[Varying $c$ alone is insufficient]
  As noted in Warning~\ref{warn:not-convolution}, varying only the scalar
  $c$ provides a one-parameter family, which is generically too small for
  Sard--Smale (the cokernel at a degenerate zero may have dimension
  $>1$).  A richer family of perturbations is needed.
\end{remark}



\section{Compactness and Finite Eigenvalue Count}
\label{sec:finite-count}

The qualitative gap around $\Lambda_0$ ultimately comes from compactness
plus finite-dimensional analytic structure.  The compactness statements
below are also the starting point for any effective version of the gap.

\subsection{Compactness Modulo Symmetry}
\label{sec:compactness}
% =======================================================================

\subsubsection{Regularity bootstrap}

\begin{proposition}[Smoothness and gauge-fixed uniform bounds]
  \label{prop:bootstrap}
  \proven\quad
  If $f\in L^2(S^1)$ satisfies $E^*(|Ef|^4Ef)=\Lambda f$ with
  $\Lambda>0$, then $f\in C^\infty(S^1)$.  Moreover, for each
  $k\geq 0$ there exists $C_k=C_k(\varepsilon)$ such that
  $\norm{f}_{C^k(S^1)}\leq C_k$ for all normalised critical points
  with $\Lambda\geq\varepsilon$ in a fixed translation gauge; in
  particular, this applies in the nonnegative slice used in
  Lemma~\ref{lem:positive-gauge}.
\end{proposition}

\begin{proof}[Proof sketch]
  Write $u=Ef$, $g=|u|^4u$.  Since $\wh u=f\,d\sigma$ is supported on
  $S^1$, the Fourier transform $\wh g$ is supported in the Minkowski sum
  $S^1+S^1+S^1+S^1+S^1=\overline{B(0,5)}$ (using $|u|^4u=u^3\bar u^2$
  and $\wh{\bar u}$ supported on $-S^1=S^1$).  Thus $g$ is
  \emph{frequency-localised} to $B(0,5)$.

  The Euler--Lagrange equation gives $f=\Lambda^{-1}\wh g\big|_{S^1}$,
  so $f$ inherits the regularity of~$\wh g$ restricted to the compact
  curve $S^1\subset B(0,5)$.  The $n$-th Fourier coefficient of $f$
  on $S^1$ is
  \begin{equation}\label{eq:fourier-bootstrap}
    \wh f(n)
    = \Lambda^{-1}\cdot 2\pi(-i)^{|n|}
      \int_0^\infty g_n(r)\,J_{|n|}(r)\,r\,dr,
  \end{equation}
  where $g_n(r)$ is the $n$-th angular Fourier coefficient of~$g$.
  For $|n|\gg 1$, $J_{|n|}(r)$ is exponentially small for $r<|n|/2$
  (Debye bounds), and the contribution from $r>|n|/2$ is controlled by
  the $L^{6/5}$ tail of~$g$ via H\"older and the Bessel asymptotics
  $J_\nu(r)=O(r^{-1/2})$ for $r>\nu$.  An inductive argument (improved
  decay of $\wh f(n)$ feeds back into improved regularity of $f$ and
  hence of $u$ and~$g$) yields $|\wh f(n)|=O(|n|^{-N})$ for all~$N$.
  The constants are obtained from the $L^2$ normalisation, the lower
  bound $\Lambda\geq\varepsilon$, and the Euler--Lagrange equation, by
  this energy/bootstrap scheme.

  The translation-gauge qualification is necessary.  If
  $f_x(\omega)=e^{-ix\cdot\omega}f(\omega)$, then $f_x$ is again a
  normalised critical point with the same~$\Lambda$, while
  $\norm{f_x}_{C^1}$ grows like~$|x|$ in general.  Once the translation
  parameter is fixed, or once one works in the nonnegative slice below,
  this obstruction is absent.
\end{proof}

\begin{corollary}[Uniform pointwise decay in a fixed gauge]
  \label{cor:uniform-decay}
  For all normalised critical points in a fixed translation gauge with
  $\Lambda\geq\varepsilon$,
  \[
    |Ef(x)|\leq C(\varepsilon)\,(1+|x|)^{-1/2}
    \qquad\text{for all }x\in\R^2.
  \]
  Equivalently, for an arbitrary translated representative the same
  estimate holds around its translation centre.
\end{corollary}
\begin{proof}
  By the bootstrap, $\norm{f}_{C^1}\leq C(\varepsilon)$.  The standard
  stationary-phase estimate for extension from a $C^\infty$ curve in
  $\R^2$ with $C^1$ density gives $|Ef(x)|=O(|x|^{-1/2})$
  for $|x|\geq 1$, and
  $|Ef(x)|\leq(2\pi)^{1/2}\norm{f}_{L^2}\leq(2\pi)^{1/2}$ everywhere.
\end{proof}

\subsubsection{Profile decomposition for the circle extension}

The extension $E:L^2(S^1)\to L^6(\R^2)$ is not compact
(Proposition~\ref{prop:E-not-compact}), and the sole obstruction is
the translation symmetry.  The following decomposition, analogous to
results of Merle--Vega and Shao for the paraboloid, captures this
defect of compactness completely.

\begin{theorem}[Profile decomposition]
  \label{thm:profile-decomp}
  \conditional\quad
  Let $f_j\in L^2(S^1)$ with $\norm{f_j}_2\leq C$.  After passing
  to a subsequence, there exist profiles $\phi^k\in L^2(S^1)$
  ($k=1,2,\ldots$) and group elements
  $g_j^k=(e^{i\alpha_j^k},R_j^k,x_j^k)\in G$ satisfying the
  \emph{divergence condition}
  $|x_j^k-x_j^l|\to\infty$ for all $k\neq l$,
  such that for each $K\geq 1$,
  \begin{equation}\label{eq:profile-decomp}
    f_j = \sum_{k=1}^K(g_j^k)^{-1}\cdot\phi^k + r_j^K,
  \end{equation}
  with the decoupling identities
  \begin{align}
    \norm{f_j}_2^2
    &= \textstyle\sum_{k=1}^K\norm{\phi^k}_2^2
      + \norm{r_j^K}_2^2 + o(1),
    \label{eq:L2-decouple}\\
    \norm{Ef_j}_6^6
    &= \textstyle\sum_{k=1}^K\norm{E\phi^k}_6^6
      + \norm{Er_j^K}_6^6 + o(1),
    \label{eq:L6-decouple}
  \end{align}
  and $\norm{Er_j^K}_{L^6}\to 0$ as $K\to\infty$ (after further
  subsequence extraction at each stage).
\end{theorem}

The $L^2$ decoupling~\eqref{eq:L2-decouple} is elementary:
for $k\neq l$,
$\ip{e^{-ix_j^k\cdot\omega}\phi^k}{e^{-ix_j^l\cdot\omega}\phi^l}
=\int_{S^1}\phi^k\overline{\phi^l}\,
e^{i(x_j^l-x_j^k)\cdot\omega}\,d\sigma\to 0$
by the Riemann--Lebesgue lemma on $S^1$.

The $L^6$ decoupling~\eqref{eq:L6-decouple} is deeper.  Since
$E(e^{-ix_j^k\cdot\omega}\phi^k)(x)=(E\phi^k)(x-x_j^k)$, distinct
profiles produce extensions centred at diverging locations.  A
Br\'ezis--Lieb argument using the Herglotz wave decay
$|E\phi^k(x)|=O(|x|^{-1/2})$ (valid once $\phi^k$ is smooth;
for $L^2$ profiles, one approximates and uses density) then gives the
$L^6$ asymptotic additivity.

The \emph{construction} of the decomposition proceeds iteratively.
The key input is an inverse Strichartz inequality for $S^1$:

\begin{lemma}[Inverse Strichartz for $S^1$]
  \label{lem:inv-strichartz}
  If $\norm{f_j}_2=1$ and $\norm{Ef_j}_{L^6}\geq\delta>0$, then
  there exist $x_j\in\R^2$ such that the re-centred sequence
  $\tilde f_j(\omega):=e^{ix_j\cdot\omega}f_j(\omega)$ has a
  subsequence converging weakly to some $\phi\neq 0$ in $L^2(S^1)$.
\end{lemma}

\begin{proof}
  Since $\norm{Ef_j}_6^6\geq\delta^6$, a covering argument gives
  $x_j\in\R^2$ with $\int_{B(x_j,1)}|Ef_j|^6\geq c\delta^6$.
  Set $\tilde u_j(x):=Ef_j(x+x_j)=E\tilde f_j(x)$.  Then
  $\int_{B(0,1)}|\tilde u_j|^6\geq c\delta^6$.

  Every Herglotz wave satisfies $(\Delta+1)u=0$ and
  $\norm{u}_{L^\infty}\leq(2\pi)^{1/2}\norm{f}_2=(2\pi)^{1/2}$.
  By interior Schauder estimates for $\Delta u+u=0$, the family
  $\{\tilde u_j\}$ is equicontinuous on $\overline{B(0,1)}$.
  Arzel\`a--Ascoli extracts a uniformly convergent subsequence
  $\tilde u_j\to u_\infty$ on $\overline{B(0,1)}$.

  Meanwhile, $\tilde f_j\rightharpoonup\phi$ (Banach--Alaoglu,
  after subsequence).  Weak convergence gives
  $E\tilde f_j(x)\to E\phi(x)$ pointwise for every~$x$
  (since $\omega\mapsto e^{ix\cdot\omega}\in L^2(S^1)$).  So
  $u_\infty=E\phi$.  From
  $\int_{B(0,1)}|u_\infty|^6\geq c\delta^6>0$, we get
  $\phi\neq 0$.
\end{proof}

\begin{remark}[Simplification from compact frequency support]
  For the paraboloid $\{(\xi,|\xi|^2)\}$, profile decompositions must
  account for parabolic rescalings (a second family of symmetries).
  For $S^1$, the frequency set is compact with no scaling symmetry,
  so translations are the sole defect of compactness.  Moreover, the
  inverse Strichartz lemma follows from elliptic regularity for
  $\Delta u+u=0$, rather than from bilinear restriction estimates.
\end{remark}

\subsubsection{Profiles of critical sequences inherit the equation}

The following result is the key structural input: if the original
sequence consists of critical points, then each profile in the
decomposition is itself a critical point.

\begin{proposition}[Each profile is a critical point]
  \label{prop:profile-critical}
  \proven\quad
  Let $(f_j,\Lambda_j)$ be normalised critical points with
  $\Lambda_j\to\Lambda_\infty>0$, and let $\phi^1,\phi^2,\ldots$
  be the profiles from Theorem~\ref{thm:profile-decomp} (re-indexed
  so that $x_j^1=0$).  Then each profile satisfies
  \begin{equation}\label{eq:profile-EL}
    E^*(|E\phi^k|^4E\phi^k) = \Lambda_\infty\phi^k.
  \end{equation}
  In particular, $\psi^k:=\phi^k/\norm{\phi^k}_2$ is a normalised
  critical point with eigenvalue
  $\widetilde\Lambda_k=\Lambda_\infty/\norm{\phi^k}_2^4$.
\end{proposition}

\begin{proof}
  We prove~\eqref{eq:profile-EL} for $k=1$; the case $k\geq 2$
  follows by applying the translation $y_j^k:=x_j^k-x_j^1$
  (which diverges) and repeating the argument.

  Set $v^l:=E\phi^l$ and $w_j:=Er_j$.  Write
  $\tilde f_j:=e^{ix_j^1\cdot\omega}f_j$ for the re-centred
  sequence, so
  $\tilde u_j:=E\tilde f_j
  =v^1+\sum_{l\geq 2}v^l(\cdot-y_j^l)+w_j$.
  The Euler--Lagrange equation, tested against arbitrary
  $h\in L^2(S^1)$, reads
  \begin{equation}\label{eq:EL-tested}
    \Lambda_j\ip{\tilde f_j}{h}_{L^2(S^1)}
    = \int_{\R^2}|\tilde u_j|^4\tilde u_j\,\overline{Eh}\,dx.
  \end{equation}

  \textbf{Left side.}
  Since $\tilde f_j=\phi^1+\sum_{l\geq 2}e^{-iy_j^l\cdot\omega}\phi^l
  +r_j$ and $\ip{e^{-iy_j^l\cdot\omega}\phi^l}{h}\to 0$ by
  Riemann--Lebesgue ($|y_j^l|\to\infty$), the left side converges to
  $\Lambda_\infty\ip{\phi^1}{h}$.

  \textbf{Right side.}
  Decompose $|\tilde u_j|^4\tilde u_j=|v^1|^4v^1+N_j$, where
  $N_j$ collects all terms containing at least one factor of
  $\eta_j:=\tilde u_j-v^1=\sum_{l\geq 2}v^l(\cdot-y_j^l)+w_j$.
  We show $\int N_j\,\overline{Eh}\to 0$.

  \emph{Step (i): $\eta_j$ vanishes locally in $L^6$.}
  For each fixed $R>0$:
  $\norm{v^l(\cdot-y_j^l)}_{L^6(B(0,R))}
  =\norm{v^l}_{L^6(B(-y_j^l,R))}\to 0$ as $j\to\infty$ (since
  $v^l\in L^6(\R^2)$ and $B(-y_j^l,R)$ escapes to infinity), and
  $\norm{w_j}_{L^6}\to 0$.  Hence
  $\norm{\eta_j}_{L^6(B(0,R))}\to 0$.

  \emph{Step (ii): individual translated-profile terms vanish.}
  Fix $l\geq 2$.  The typical cross term to estimate is
  $I_j^l:=\int|v^1(x)|^4\,v^l(x-y_j^l)\,\overline{Eh(x)}\,dx$.
  Split $\R^2$ into $B(0,R)$ and its complement.
  \begin{itemize}
    \item On $B(0,R)$: $|v^1|\leq(2\pi)^{1/2}$, so
      $|I_j^l|_{B(0,R)}|\leq C\norm{v^l(\cdot-y_j^l)}_{L^6(B(0,R))}
      \norm{Eh}_{L^6}\to 0$ by step~(i).
    \item On $\R^2\setminus B(0,R)$:
      $|I_j^l|_{\R^2\setminus B}|\leq C\norm{v^l}_6
      \norm{Eh}_{L^6(|x|>R)}\to 0$ as $R\to\infty$ (since
      $Eh\in L^6$ has vanishing tails).
  \end{itemize}
  An $\varepsilon/3$ argument in $R$ and $j$ gives $I_j^l\to 0$.

  \emph{Step (iii): the remainder.}
  Terms involving $w_j=Er_j$ are bounded by
  $C\norm{w_j}_6^{}\norm{Eh}_6\to 0$ (using
  $\norm{v^1}_\infty\leq C$ and H\"older).

  \emph{Step (iv): higher-order cross terms.}
  Every monomial in $N_j$ has at least one factor $\eta_j$ or
  $\bar\eta_j$.  The same local-vanishing argument applies: on
  $B(0,R)$, the factor of $\eta_j$ is small in $L^6$; on the
  complement, $Eh$ is small in $L^6$.

  \medskip
  Combining: the right side of~\eqref{eq:EL-tested} converges to
  $\int|v^1|^4v^1\,\overline{Eh}
  =\ip{E^*(|E\phi^1|^4E\phi^1)}{h}$.
  Since $h$ is arbitrary,
  $E^*(|E\phi^1|^4E\phi^1)=\Lambda_\infty\phi^1$.

  The rescaling: if $\alpha_k:=\norm{\phi^k}_2>0$, then
  $\psi^k:=\phi^k/\alpha_k$ satisfies
  $E^*(|E\psi^k|^4E\psi^k)
  =\alpha_k^{-5}E^*(|E\phi^k|^4E\phi^k)
  =\Lambda_\infty\alpha_k^{-5}\phi^k
  =(\Lambda_\infty/\alpha_k^4)\psi^k$.
\end{proof}

\subsubsection{Existence of a maximiser}

\begin{theorem}[Existence of extremisers for $S^1$ restriction]
  \label{thm:maximiser}
  \proven\quad
  The sharp restriction constant $\mathcal R$ is attained: there exists
  $f_*\in L^2(S^1)$ with $\norm{f_*}_2=1$ and
  $\norm{Ef_*}_{L^6}=\mathcal R$.
\end{theorem}

\begin{proof}
  Let $f_j$ be a maximising sequence: $\norm{f_j}_2=1$,
  $\Lambda_j:=\norm{Ef_j}_6^6\to\mathcal R^6>0$.  Apply the profile
  decomposition (Theorem~\ref{thm:profile-decomp}).

  \textbf{Step 1: at least one profile.}
  Since $\norm{Er_j^K}_6\to 0$ as $K\to\infty$ while
  $\Lambda_j\to\mathcal R^6>0$, the
  decoupling~\eqref{eq:L6-decouple} forces $\phi^1\neq 0$.

  \textbf{Step 2: no splitting.}
  By the restriction inequality, $\norm{E\phi^k}_6^6
  \leq\mathcal R^6\norm{\phi^k}_2^6$.  Set $\alpha_k:=\norm{\phi^k}_2$
  and $a_k:=\alpha_k^2$.  The decoupling gives
  $\mathcal R^6\leq\mathcal R^6\sum_{k}a_k^3$.
  From~\eqref{eq:L2-decouple}, $\sum_k a_k\leq 1$.

  If $M\geq 2$ profiles are nonzero, all $a_k>0$ and $\sum a_k\leq 1$.
  But the multinomial expansion
  $\bigl(\sum_k a_k\bigr)^3
  =\sum_k a_k^3+3\sum_{k\neq l}a_k^2a_l+6\sum_{k<l<m}a_ka_la_m$
  shows $\sum_k a_k^3<\bigl(\sum_k a_k\bigr)^3\leq 1$ (since the
  cross terms are strictly positive).  This gives
  $\mathcal R^6<\mathcal R^6$, a contradiction.

  \textbf{Step 3: convergence.}
  With a single profile: $f_j=g_j^{-1}\cdot\phi^1+r_j$,
  $\norm{r_j}_2\to 0$, $\norm{\phi^1}_2=1$.  Hence
  $g_j\cdot f_j\to\phi^1$ strongly, and $\phi^1$ is a maximiser.
\end{proof}

\begin{remark}[The homogeneity gap]
  The argument uses $6>2$: the $L^6$ norm is superquadratic in~$f$,
  making $\sum a_k^3<(\sum a_k)^3$ when there are multiple profiles.
  This mechanism is specific to maximising sequences and does not
  extend to sub-maximal critical points.
\end{remark}

\subsubsection{The splitting obstruction for general critical sequences}

For a general sequence of critical points with
$\Lambda_j\to\Lambda_\infty\in[\varepsilon,\mathcal R^6)$, the
strict-convexity argument breaks down: multi-profile decompositions
are \emph{algebraically consistent} with the Euler--Lagrange equation.

\begin{proposition}[Structure of multi-profile decompositions at
  critical sequences]
  \label{prop:multi-profile}
  \proven\quad
  Suppose a sequence of normalised critical points with
  $\Lambda_j\to\Lambda_\infty>0$ admits a profile decomposition
  with $M\geq 2$ nonzero profiles $\phi^1,\ldots,\phi^M$.  Then:
  \begin{enumerate}[leftmargin=2em]
    \item Each normalised profile $\psi^k=\phi^k/\alpha_k$ (where
      $\alpha_k:=\norm{\phi^k}_2$) is a critical point with
      eigenvalue
      $\widetilde\Lambda_k=\Lambda_\infty/\alpha_k^4
      >\Lambda_\infty$.
    \item The number of profiles satisfies
      $M\leq(\mathcal R^6/\varepsilon)^{1/2}$.
    \item The identity $\sum_k\alpha_k^6\widetilde\Lambda_k
      =\Lambda_\infty\sum_k\alpha_k^2=\Lambda_\infty$ holds
      automatically, so the energy decomposition produces no
      contradiction.
  \end{enumerate}
\end{proposition}

\begin{proof}
  (1) Proposition~\ref{prop:profile-critical} gives
  $E^*(|E\phi^k|^4E\phi^k)=\Lambda_\infty\phi^k$.  Rescaling:
  $E^*(|E\psi^k|^4E\psi^k)=(\Lambda_\infty/\alpha_k^4)\psi^k$,
  so $\widetilde\Lambda_k=\Lambda_\infty/\alpha_k^4$.  Since
  $\alpha_k^2<\sum_l\alpha_l^2=1$, we have $\alpha_k<1$ and
  $\widetilde\Lambda_k>\Lambda_\infty$.

  (2) From $\widetilde\Lambda_k\leq\mathcal R^6$:
  $\alpha_k^4\geq\Lambda_\infty/\mathcal R^6
  \geq\varepsilon/\mathcal R^6$, so
  $\alpha_k^2\geq(\varepsilon/\mathcal R^6)^{1/2}$.  The constraint
  $\sum_k\alpha_k^2=1$ gives the bound on~$M$.

  (3) $\sum_k\alpha_k^6\widetilde\Lambda_k
  =\sum_k\alpha_k^6\cdot\Lambda_\infty/\alpha_k^4
  =\Lambda_\infty\sum_k\alpha_k^2=\Lambda_\infty$. \qedhere
\end{proof}

Part~(3) is the central difficulty:
\textbf{the Euler--Lagrange equation alone does not prevent a critical
sequence from splitting into bubbles at diverging locations, each
individually a critical point at a higher eigenvalue level.}
This is the open content of Theorem~\ref{thm:compactness-mod-G}:

\begin{theorem}[Compactness modulo $G$]
  \label{thm:compactness-mod-G}
  \open\quad
  Let $(f_j,\Lambda_j)$ be normalised critical points with
  $\Lambda_j\geq\varepsilon>0$.  Then there exist $g_j\in G$ and a
  subsequence such that $g_j\cdot f_j\to f_\infty$ strongly in
  $L^2(S^1)$ and $\Lambda_j\to\Lambda_\infty=\norm{Ef_\infty}_6^6$.
\end{theorem}

Equivalently: every critical sequence has \textbf{exactly one
nonzero profile} in its decomposition.

\subsubsection{Approaches to rule out splitting}

We describe three general strategies that may close the gap, and then an
extremiser-specific refinement using nonnegative representatives.

\paragraph{(A) Descending induction on critical values.}
At the maximal level $\Lambda=\mathcal R^6$,
Theorem~\ref{thm:maximiser} proves single-profile concentration.
Suppose inductively that compactness modulo~$G$ holds for critical
sequences with $\Lambda_j\to\Lambda'$ for all
$\Lambda'>\Lambda_\infty$.  A splitting at level~$\Lambda_\infty$
produces profiles at levels $\widetilde\Lambda_k>\Lambda_\infty$,
where the inductive hypothesis already gives finiteness of shapes.
The splitting identity $\sum_k\alpha_k^2=1$,
$\widetilde\Lambda_k=\Lambda_\infty/\alpha_k^4$ imposes algebraic
constraints relating the profile masses to the (now finitely many)
critical values above $\Lambda_\infty$.

\emph{Difficulty}: the base case gives one shape (the maximiser);
but the inductive step requires both finiteness and nondegeneracy at
every intermediate level, creating a dependency between this theorem
and the results of Section~\ref{sec:nondeg}.  Breaking this
circularity requires either a direct nondegeneracy argument at all
levels, or an independent route to single-profile concentration.

\paragraph{(B) Pohozaev--Morawetz monotonicity.}
A localised identity for $u=c(|u|^4u\ast\wh\sigma)$ could yield
a monotonicity formula controlling the spatial distribution of
$|u|^6$ and preventing mass separation.  The starting point would be
to multiply the equation by a radial weight $a(|x|)\,\bar u$ and
integrate; the resulting identity relates $\int a'|u|^6$ to boundary
terms and the nonlocal convolution structure.

\emph{Difficulty}: the convolution against $\wh\sigma$ (which decays
only like $|x|^{-1/2}$) introduces long-range interactions that
resist clean localisation.

\paragraph{(C) Unique continuation for the integrated equation.}
If one can show that two critical points $\phi^1,\phi^2$ at levels
$\widetilde\Lambda_1,\widetilde\Lambda_2>\Lambda_\infty$ cannot
coexist as profiles of a common critical sequence---i.e., the
relation
\[
  f_j\approx e^{-ix_j^1\cdot\omega}\phi^1+e^{-ix_j^2\cdot\omega}\phi^2,
  \qquad |x_j^1-x_j^2|\to\infty,
\]
contradicts $E^*(|Ef_j|^4Ef_j)=\Lambda_jf_j$ at order beyond the
profile-level Euler--Lagrange---this would rule out splitting.  The
relevant obstruction would arise from cross terms at intermediate
spatial scales $|x|\sim|x_j^1-x_j^2|^{1/2}$ where the two Herglotz
waves overlap at polynomial (not exponential) level.

\paragraph{(D) Positive gauge for the extremiser search.}
For the sharp constant, Lemma~\ref{lem:positive-reduction} allows the
search to be restricted to nonnegative critical points.  This is not
available for arbitrary critical sequences, but it removes much of the
translation ambiguity: a nonnegative representative cannot be translated
nontrivially and remain nonnegative.

\begin{lemma}[Nonnegativity fixes the translation parameter]
  \label{lem:positive-gauge}
  Let $0\neq f\in L^2(S^1)$ with $f\geq 0$ a.e.  If
  \[
    e^{i\alpha}e^{-ix_0\cdot\omega}f(R^{-1}\omega)\geq 0
    \qquad\text{for a.e. }\omega\in S^1,
  \]
  then $x_0=0$.  In that case the remaining phase must satisfy
  $e^{i\alpha}=1$ on the support of $f$.
\end{lemma}

\begin{proof}
  The set where $f(R^{-1}\omega)>0$ has positive measure.  On this set,
  the phase $e^{i\alpha}e^{-ix_0\cdot\omega}$ must equal~$1$.  If
  $x_0\neq 0$, the set
  $\{\omega\in S^1:e^{i\alpha}e^{-ix_0\cdot\omega}=1\}$ is finite, since
  it is a finite union of level sets $x_0\cdot\omega=\alpha+2\pi m$.
  This cannot contain a positive-measure subset of~$S^1$.  Hence $x_0=0$.
\end{proof}

The profile-decomposition consequence is the following.  By the
gauge-fixed bootstrap in Proposition~\ref{prop:bootstrap}, a
nonnegative critical sequence with $\Lambda_j\geq\varepsilon$ satisfies
$\norm{f_j}_{C^1(S^1)}\leq C_\varepsilon$.  Thus every escaping profile
vanishes: for $|x_j|\to\infty$ and $h\in C^1(S^1)$, stationary phase gives
\[
  \int_{S^1} e^{ix_j\cdot\omega} f_j(\omega)\overline{h(\omega)}
  \,d\sigma(\omega)
  = O_{C_\varepsilon,h}(|x_j|^{-1/2}),
\]
and density extends this weak convergence to all $h\in L^2(S^1)$.  In the
iterative profile construction, after subtracting finitely many previously
selected profiles, the terms involving
$e^{i(x_j^k-x_j^\ell)\cdot\omega}\phi^\ell$ also vanish by the
Riemann--Lebesgue lemma whenever the translation parameters diverge.
Thus a candidate profile with $|x_j^k|\to\infty$ has zero weak limit.

Thus the extra input needed for the extremiser search is not a new
profile-decomposition mechanism, but the standard gauge-fixed bootstrap:
the $L^2$ normalisation, the equation, and the lower bound on~$\Lambda$
give the uniform norms, while positivity supplies the gauge because pure
translation orbits leave the nonnegative slice by
Lemma~\ref{lem:positive-gauge}.  This closes the escaping-profile
obstruction for the extremiser search without proving compactness of the
full critical set.

\begin{remark}[Sufficiency for identifying $\mathcal R$]
  \label{rem:sufficiency-max}
  For the purpose of \emph{computing} the sharp constant~$\mathcal R$,
  the full Theorem~\ref{thm:compactness-mod-G} is not needed.
  Theorem~\ref{thm:maximiser} gives existence of a maximiser, and if
  the maximiser is unique modulo~$G$ (e.g., the constant function
  $f_0=(2\pi)^{-1/2}$), then $\mathcal R=\norm{Ef_0}_{L^6}$ and
  the problem is solved.  The open compactness theorem would
  additionally yield finiteness of \emph{all} critical values above
  $\varepsilon$, which is a stronger structural result.
\end{remark}

\begin{remark}[Accumulation of critical values at zero]
  \label{rem:accumulation}
  Whether critical values can accumulate at~$0$ is unknown.  The
  compactness + nondegeneracy programme gives finiteness only
  above a fixed threshold.  For the sharp constant, it suffices to
  work above any $\varepsilon$ smaller than the maximum critical
  value.
\end{remark}


\subsection{Finiteness after the positivity reduction}

\begin{theorem}[Finiteness of positive critical values in compact windows]
  \label{thm:positive-values-finite}
  \proven\quad
  Fix $0<\Lambda_-<\Lambda_+<\infty$.  Then the set
  \[
    \mathcal V_+(\Lambda_-,\Lambda_+)
    :=
    \left\{
      \Lambda\in[\Lambda_-,\Lambda_+]:
      \begin{array}{l}
      \exists f\in L^2(S^1),\ f\geq 0,\ \norm{f}_2=1,\\
      E^*(|Ef|^4Ef)=\Lambda f
      \end{array}
    \right\}
  \]
  is finite.
\end{theorem}

\begin{proof}
  Let
  \[
    \mathcal C_+
    :=
    \left\{
      (f,\Lambda):
      f\geq0,\ \norm{f}_2=1,\ \Lambda\in[\Lambda_-,\Lambda_+],\
      E^*(|Ef|^4Ef)=\Lambda f
    \right\}.
  \]

  \emph{Step 1: compactness of the positive slice.}
  By the gauge-fixed bootstrap, Proposition~\ref{prop:bootstrap}, for
  every $m\geq0$ there is a constant $C_m=C_m(\Lambda_-)$ such that
  $\norm{f}_{C^m(S^1)}\leq C_m$ for every $(f,\Lambda)\in\mathcal C_+$.
  The reason this applies without an additional translation choice is
  precisely Lemma~\ref{lem:positive-gauge}: the nonnegative slice fixes
  the noncompact translation parameter.

  Given a sequence $(f_j,\Lambda_j)\in\mathcal C_+$, pass to a subsequence
  with $\Lambda_j\to\Lambda_\infty\in[\Lambda_-,\Lambda_+]$.  The uniform
  $C^{m+1}$ bound and Arzel\`a--Ascoli give, after passing to a further
  subsequence, $f_j\to f_\infty$ in $C^m(S^1)$ for each fixed~$m$ (and in
  particular in $L^2$).  Then $f_\infty\geq0$ and
  $\norm{f_\infty}_2=1$.  Since
  \[
    T(f):=E^*(|Ef|^4Ef)
  \]
  is continuous $L^2(S^1)\to L^2(S^1)$ by Stein--Tomas and H\"older,
  passing to the limit in $T(f_j)=\Lambda_j f_j$ gives
  $T(f_\infty)=\Lambda_\infty f_\infty$.  Hence
  $(f_\infty,\Lambda_\infty)\in\mathcal C_+$.  Thus $\mathcal C_+$ is
  compact in $L^2(S^1)\times\R$.

  \emph{Step 2: local finite-dimensional analytic structure.}
  Work on the real Hilbert space $H:=L^2(S^1;\C)$ with real inner product
  $\re\ip{\cdot}{\cdot}$.  Define
  \[
    \Phi:H\times\R\to H\times\R,\qquad
    \Phi(f,\Lambda):=\bigl(T(f)-\Lambda f,\ \norm{f}_2^2-1\bigr).
  \]
  The map $T$ is a real-analytic quintic polynomial map:
  $T(f)=E^*((Ef)^3(\overline{Ef})^2)$.  At a zero $(f_0,\Lambda_0)$,
  \[
    D\Phi_{(f_0,\Lambda_0)}(h,\mu)
    =
    \left(
      (K_{u_0}-\Lambda_0 I)h-\mu f_0,\
      2\re\ip{f_0}{h}
    \right),
    \qquad u_0=Ef_0.
  \]
  By Corollary~\ref{cor:fredholm}, $K_{u_0}-\Lambda_0 I$ is Fredholm of
  index~$0$.  The displayed derivative is a finite-rank perturbation of
  $(h,\mu)\mapsto((K_{u_0}-\Lambda_0 I)h,\mu)$, and is therefore Fredholm
  of index~$0$.

  The analytic Fredholm Lyapunov--Schmidt reduction now applies: near
  each zero of~$\Phi$, the zero set $\Phi^{-1}(0)$ is homeomorphic to the
  zero set of a finite-dimensional real-analytic map.  By the standard
  local structure theorem for real-analytic sets (equivalently,
  Lojasiewicz local finiteness), this local zero set has only finitely
  many connected components, and each such component is locally
  arc-connected through finitely many analytic arcs.

  \emph{Step 3: the value is constant on local critical components.}
  Let $J(f):=\norm{Ef}_{L^6}^6$.  On $\Phi^{-1}(0)$ one has
  $\Lambda=J(f)$, because
  \[
    \Lambda\norm{f}_2^2
    =\re\ip{T(f)}{f}
    =\int_{\R^2}|Ef(x)|^6\,dx.
  \]
  If $t\mapsto(f(t),\Lambda(t))$ is a $C^1$ arc in $\Phi^{-1}(0)$, then
  $\norm{f(t)}_2=1$ and $T(f(t))=\Lambda(t)f(t)$, so
  \[
    \frac{d}{dt}J(f(t))
    =6\re\ip{T(f(t))}{f'(t)}
    =6\Lambda(t)\re\ip{f(t)}{f'(t)}
    =0.
  \]
  Therefore $J(f(t))$, and hence $\Lambda(t)$, is constant on every such
  arc.  Since the local analytic components are connected by finite
  chains of these arcs, $\Lambda$ is constant on each local connected
  component of $\Phi^{-1}(0)$.

  Finally, compactness of $\mathcal C_+$ gives a finite cover by
  neighbourhoods of the type described in Step~2.  Each neighbourhood
  contains only finitely many local critical components, and Step~3 gives
  at most one value of~$\Lambda$ on each of them.  Hence the projection
  of $\mathcal C_+$ onto the $\Lambda$ coordinate, namely
  $\mathcal V_+(\Lambda_-,\Lambda_+)$, is finite.
\end{proof}



\section{Effective Spectral Gap Around the Constant Level}
\label{sec:effective-gap}

Let
\[
  f_0=(2\pi)^{-1/2},\qquad
  \Lambda_0:=\norm{Ef_0}_{L^6}^6
  =16\pi^4\int_0^\infty J_0(r)^6\,r\,dr.
\]
If $C_{\mathrm{ST}}$ denotes any explicit Stein--Tomas bound for
$E:L^2(S^1)\to L^6(\R^2)$, then every normalised critical value satisfies
$\Lambda\leq C_{\mathrm{ST}}^6$.  Thus the natural positive window around
the constant value is
\[
  W_0:=\left[\frac12\Lambda_0,\ C_{\mathrm{ST}}^6\right].
\]
By Theorem~\ref{thm:positive-values-finite}, the set
$\mathcal V_+(W_0)$ of positive critical values in this window is finite
and contains~$\Lambda_0$.  Hence
\begin{equation}\label{eq:positive-gap-def}
  \gamma_0
  :=
  \operatorname{dist}\bigl(
    \Lambda_0,\,
    \mathcal V_+(W_0)\setminus\{\Lambda_0\}
  \bigr)
\end{equation}
is strictly positive, with the convention $\gamma_0=+\infty$ if
$\mathcal V_+(W_0)=\{\Lambda_0\}$.

The qualitative proof gives no usable lower bound on~$\gamma_0$.  The
central task is therefore to make the following implication effective:
\[
  T(f)=\Lambda f,\quad \norm f_2=1,\quad
  \abs{\Lambda-\Lambda_0}\ll1
  \quad\Longrightarrow\quad
  \operatorname{dist}(f,G\cdot f_0)\ll1.
\]
Once a solution is forced into a sufficiently small tube around the
constant orbit, quantitative nondegeneracy at~$f_0$ excludes it unless it
is the constant itself.

\subsection{Quantitative local isolation at the constant}

On the slice orthogonal to phase and translations, define the linear
nondegeneracy gap
\[
  \nu_0:=
  \min\left\{
    \frac45\Lambda_0,\;
    \inf_{n\geq2}\abs{\mathcal I_n-\Lambda_0},\;
    \inf_{n\geq2}\abs{5\mathcal I_n-\Lambda_0}
  \right\},
  \qquad
  \mathcal I_n=16\pi^4\mathcal J_n.
\]
The constant-solution nondegeneracy theorem gives $\nu_0>0$.  For an
effective radius one needs explicit constants:
\[
  \nu_0\geq\underline\nu>0,
  \qquad
  \norm{D\widetilde F(f)-D\widetilde F(g)}_{X\to X}
  \leq M_0\norm{f-g}_2
\]
on a gauge-fixed slice $X$ near~$f_0$, where
$\widetilde F$ is the sliced Euler--Lagrange map.  The standard
Newton--Kantorovich estimate then gives the explicit uniqueness radius
\begin{equation}\label{eq:nk-radius}
  r_0=\frac{2\underline\nu}{M_0}.
\end{equation}
Indeed, if $\widetilde F(f)=0=\widetilde F(f_0)$ and
$\norm{f-f_0}_2<r_0$, then
\[
  \underline\nu\norm{f-f_0}_2
  \leq
  \frac{M_0}{2}\norm{f-f_0}_2^2,
\]
which is impossible unless $f=f_0$.

Thus a quantitative spectral gap follows from any effective compactness
statement which places all $\Lambda$-nearby solutions inside this ball.
The remaining subsections describe two analytic ways to obtain such a
statement.

\subsection{Effective Fredholm-atlas route}
\label{sec:effective-fredholm-atlas}

Let
\[
  \mathcal S_\Lambda
  :=
  \{f\geq0:\ \norm f_2=1,\ T(f)=\Lambda f\}
\]
in the positive gauge.  The desired effective compactness statement is a
modulus
\begin{equation}\label{eq:effective-upper-semicontinuity}
  \operatorname{dist}\bigl(\mathcal S_\Lambda,\mathcal S_{\Lambda_0}\bigr)
  \leq \rho(\abs{\Lambda-\Lambda_0}),
  \qquad \rho(t)\to0\quad(t\downarrow0),
\end{equation}
with explicit control of~$\rho$.  If, in addition,
\[
  \mathcal S_{\Lambda_0}/G=\{[f_0]\},
\]
then choosing $\eta>0$ so that $\rho(\eta)<r_0$ proves
$\gamma_0\geq\eta$.

One way to make \eqref{eq:effective-upper-semicontinuity} quantitative is
to revisit the proof of finite eigenvalue count with constants.  Fix a
small window
\[
  I_\eta=[\Lambda_0-\eta,\Lambda_0+\eta].
\]
The ingredients needed are the following.

\begin{enumerate}[leftmargin=2em]
  \item \textbf{Effective analytic a priori bounds.}  For every solution
  with $\Lambda\in I_\eta$, prove explicit estimates in a fixed analytic
  Banach space, for example
  \[
    \norm{f}_{\mathcal A_\rho}\leq A(\eta),
  \]
  where $\mathcal A_\rho$ denotes functions holomorphic in a complex
  strip $\{|\operatorname{Im}\theta|<\rho\}$ with a Banach norm such as
  the sup norm on the strip.  The point is not to truncate Fourier
  series, but to put the whole solution set in one compact analytic
  class.  The inclusion $\mathcal A_\rho\hookrightarrow \mathcal A_{\rho'}$
  is compact for $\rho'<\rho$, and the constants $A(\eta),\rho,\rho'$
  give a computable compactness modulus for the possible shapes.

  \item \textbf{Effective finite Fredholm atlas at level $\Lambda_0$.}
  Cover $\mathcal S_{\Lambda_0}$ by finitely many gauge-fixed charts.
  At a chart centre $f_*$, write the constrained equation as
  \[
    F(f,\Lambda)=0,
    \qquad
    F(f,\Lambda):=T(f)-\Lambda f,
    \qquad \norm f_2=1,
  \]
  on a slice.  Since $D_fF(f_*,\Lambda_0)$ is Fredholm, split the slice as
  \[
    X=K_*\oplus X_*,\qquad
    K_*:=\ker D_fF(f_*,\Lambda_0),
  \]
  and solve the $X_*$ component by a quantitative inverse function
  theorem.  This reduces the local zero set to a finite-dimensional
  analytic equation
  \[
    \Phi_*(a,\Lambda)=0,
    \qquad a\in K_*.
  \]
  The number of active parameters is therefore bounded by
  $\dim K_*$, which is finite.  An effective proof needs an explicit upper
  bound for these kernel dimensions, or a direct computable basis for the
  kernels in the relevant atlas.

  \item \textbf{Quantitative inverse bounds off the singular set.}  Away
  from the singular locus
  \[
    \Sigma_{\Lambda_0}:=
    \{f\in\mathcal S_{\Lambda_0}:D_fF(f,\Lambda_0)
      \text{ is not invertible on the slice}\},
  \]
  one needs a lower bound
  \begin{equation}\label{eq:away-singular-inverse}
    \norm{D_fF(f,\Lambda_0)h}_2
    \geq m(\delta)\norm h_2
  \end{equation}
  whenever $\operatorname{dist}(f,\Sigma_{\Lambda_0})\geq\delta$.
  Together with a Lipschitz bound for $D_fF$, this gives a quantitative
  local inverse theorem: solution sets at nearby $\Lambda$ move by at most
  $O(\abs{\Lambda-\Lambda_0}/m(\delta))$ in these regular charts.

  \item \textbf{Effective treatment of singular charts.}  Near
  $\Sigma_{\Lambda_0}$, the Lyapunov--Schmidt maps
  $\Phi_*(a,\Lambda)$ are finite-dimensional real-analytic systems.  A
  quantitative Lojasiewicz or Weierstrass-preparation estimate would bound
  how branches at $\Lambda\neq\Lambda_0$ can leave the singular set.  If
  the only singular directions are symmetry directions and the slice has
  removed them, this step collapses to the ordinary inverse theorem.
\end{enumerate}

This is the precise form of the ``skewering'' strategy.  The a priori
bounds make the window compact in an analytic topology with a computable
modulus; the Fredholm atlas reduces the infinite-dimensional problem to
finitely many local finite-dimensional systems; inverse estimates away
from the singular set control the motion of solution sets as~$\Lambda$
varies; and the singular charts are handled by finite-dimensional analytic
estimates.  The output is an explicit function~$\rho$ in
\eqref{eq:effective-upper-semicontinuity}.  Combining
\eqref{eq:effective-upper-semicontinuity} with the local radius
\eqref{eq:nk-radius} gives an effective nonlinear spectral gap.

A stronger but cleaner sufficient input is the following stability
inequality:
\begin{equation}\label{eq:global-stability-main}
  \Lambda_0-J(f)
  \geq c_*\operatorname{dist}(f,G\cdot f_0)^2,
  \qquad
  J(f):=\norm{Ef}_6^6,
\end{equation}
for all nonnegative normalised $f$ in the relevant window.  This is a
Bianchi--Egnell type statement for the sharp restriction inequality.  If
\eqref{eq:global-stability-main} is available, the Fredholm-atlas
machinery is unnecessary for the sharp constant: any nonlinear eigenpair
with $\Lambda$ close to~$\Lambda_0$ is forced directly into the local
nondegeneracy ball.

\subsection{What has to be made effective in one chart}
\label{sec:effective-one-chart}

We spell out the local reduction because this is the point at which the
infinite-dimensional problem becomes a finite-dimensional question about
zero sets.  Work in a gauge-fixed analytic Banach space $X$ for the
unknown and an analytic target space $Y$ for the residual.  At a solution
$(f_*,\Lambda_0)$ set
\[
  F(f,\Lambda):=T(f)-\Lambda f,\qquad
  L_*:=D_fF(f_*,\Lambda_0).
\]
After imposing the normalisation and gauge conditions, $L_*:X\to Y$ is
Fredholm.  Choose splittings
\[
  X=K_*\oplus X_*,
  \qquad K_*=\ker L_*,
  \qquad
  Y=Y_*\oplus C_*,
\]
with $C_*$ finite-dimensional and $L_*:X_*\to Y_*$ an isomorphism.  Let
$Q_*:Y\to Y_*$ and $\Pi_*:Y\to C_*$ be the corresponding projections.
Writing
\[
  f=f_*+a+w,\qquad a\in K_*,\quad w\in X_*,
\]
the equation splits into
\[
  Q_*F(f_*+a+w,\Lambda)=0,
  \qquad
  \Pi_*F(f_*+a+w,\Lambda)=0.
\]
If
\[
  B_*:=\norm{(Q_*L_*|_{X_*})^{-1}}_{Y_*\to X_*}
\]
is explicit and the derivative of
$Q_*F$ is Lipschitz with constant $M_*$ on a ball of radius $r_*$, then
the quantitative inverse function theorem solves the first equation as
\[
  w=\Psi_*(a,\Lambda)
\]
for
$\norm{a}_{K_*}+|\Lambda-\Lambda_0|\leq c/(B_*M_*)$.  The remaining
equation is finite-dimensional:
\begin{equation}\label{eq:LS-reduced-map}
  \Phi_*(a,\Lambda)
  :=
  \Pi_*F(f_*+a+\Psi_*(a,\Lambda),\Lambda)=0,
  \qquad a\in K_*.
\end{equation}

Thus the constants needed from the Fredholm reduction are concrete:
\[
  \dim K_*,
  \qquad
  B_*,
  \qquad
  M_*,
  \qquad
  r_*,
  \qquad
  \norm{\partial_\Lambda\Phi_*}.
\]
Away from singular points of $\Phi_*(\cdot,\Lambda_0)$, a lower bound on
the smallest singular value of $D_a\Phi_*$ gives the local motion of zero
sets by the ordinary inverse theorem.  At singular points, the required
replacement is an effective finite-dimensional analytic estimate, such as
a Lojasiewicz distance inequality or a Weierstrass-preparation bound for
\(\Phi_*\).  This is now a problem about finitely many analytic maps on
compact subsets of Euclidean space, not a problem about global Fourier
tails.

\subsection{Quantifying the finite-dimensional zero-set constants}
\label{sec:quant-lojasiewicz-chart}

We now isolate the finite-dimensional problem produced by the previous
subsection.  Fix one Lyapunov--Schmidt chart and write
\[
  \Phi(a,\lambda)=0,
  \qquad
  a\in U\subset\R^d,\qquad \lambda\in I_\eta.
\]
Let
\[
  Z_0:=\{a\in K:\Phi(a,\Lambda_0)=0\},
  \qquad K\Subset U
\]
be the part of the level-$\Lambda_0$ zero set covered by the chart.  To
control nearby levels, it is enough to prove an effective distance
inequality
\begin{equation}\label{eq:chart-loj}
  \operatorname{dist}(a,Z_0)
  \leq C_{\mathrm{Loj}}\,
  \norm{\Phi(a,\Lambda_0)}^{\theta}
  \qquad (a\in K),
\end{equation}
together with a bound
\begin{equation}\label{eq:lambda-lip-Phi}
  \norm{\Phi(a,\lambda)-\Phi(a,\Lambda_0)}
  \leq L_\lambda |\lambda-\Lambda_0|
  \qquad (a\in K,\ \lambda\in I_\eta).
\end{equation}
Then every zero \(a_\lambda\in K\) of \(\Phi(\cdot,\lambda)\) satisfies
\begin{equation}\label{eq:chart-zero-motion}
  \operatorname{dist}(a_\lambda,Z_0)
  \leq
  C_{\mathrm{Loj}}L_\lambda^\theta
  |\lambda-\Lambda_0|^\theta .
\end{equation}
This is the local quantitative upper-semicontinuity estimate.  The lift
from \(a\)-space back to the function space only costs the Lipschitz
constant of the chart map
\[
  a,\lambda\mapsto f_*+a+\Psi(a,\lambda).
\]

\begin{proposition}[One-chart spectral-gap modulus]
  \label{prop:one-chart-gap-modulus}
  Suppose a chart satisfies \eqref{eq:chart-loj} and
  \eqref{eq:lambda-lip-Phi}.  Suppose also that the chart map satisfies
  \[
    \norm{\Psi(a,\lambda)-\Psi(b,\Lambda_0)}_X
    \leq L_a\norm{a-b}+L_\lambda'|\lambda-\Lambda_0|
  \]
  on the chart.  If the level-$\Lambda_0$ zero set in the chart is
  contained in the constant orbit and the local exclusion radius around
  that orbit is \(r_0\), then this chart contains no nonconstant solution
  with \(0<|\lambda-\Lambda_0|<\eta_*\), where any \(\eta_*>0\)
  satisfying
  \begin{equation}\label{eq:eta-chart-condition}
    L_a C_{\mathrm{Loj}}L_\lambda^\theta\eta_*^\theta
    +L_\lambda'\eta_*<r_0
  \end{equation}
  is admissible.
\end{proposition}

\begin{proof}
  Let \(a_\lambda\) be a zero of \(\Phi(\cdot,\lambda)\).  By
  \eqref{eq:chart-zero-motion}, choose \(a_0\in Z_0\) with
  \[
    \norm{a_\lambda-a_0}
    \leq C_{\mathrm{Loj}}L_\lambda^\theta
    |\lambda-\Lambda_0|^\theta .
  \]
  The corresponding functions satisfy
  \[
    \norm{f(a_\lambda,\lambda)-f(a_0,\Lambda_0)}_X
    \leq
    L_a C_{\mathrm{Loj}}L_\lambda^\theta
    |\lambda-\Lambda_0|^\theta
    +L_\lambda'|\lambda-\Lambda_0|.
  \]
  If \(|\lambda-\Lambda_0|<\eta_*\), condition
  \eqref{eq:eta-chart-condition} places the solution inside the local
  exclusion ball around the level-\(\Lambda_0\) zero set.  Local
  nondegeneracy then forces the solution to be the constant one, hence
  \(\lambda=\Lambda_0\).
\end{proof}

Thus the problem is reduced to obtaining explicit
\(C_{\mathrm{Loj}}\) and \(\theta\) for each finite-dimensional
analytic map \(\Phi(\cdot,\Lambda_0)\).  There are two regimes.

\paragraph{Regular charts.}
If
\[
  s_*:=
  \inf_{a\in K\cap Z_0}
  \sigma_{\min}\bigl(D_a\Phi(a,\Lambda_0)|_{N_aZ_0}\bigr)>0,
\]
where \(N_aZ_0\) is a chosen normal complement to the smooth zero set,
then \(\theta=1\).  If \(D_a\Phi\) is \(M_\Phi\)-Lipschitz on a tubular
neighbourhood of radius \(r\) around \(Z_0\), the quantitative inverse
theorem gives
\[
  \operatorname{dist}(a,Z_0)
  \leq 2s_*^{-1}\norm{\Phi(a,\Lambda_0)}
\]
whenever \(\operatorname{dist}(a,Z_0)\leq s_*/M_\Phi\).  Outside this
tube, compactness gives a positive residual floor
\[
  m_{\mathrm{far}}:=
  \inf_{\substack{a\in K\\
  \operatorname{dist}(a,Z_0)\geq s_*/M_\Phi}}
  \norm{\Phi(a,\Lambda_0)}.
\]
Combining the two regions gives an explicit global constant on \(K\):
\[
  C_{\mathrm{Loj}}
  =
  \max\left\{
    2s_*^{-1},\
    \frac{\operatorname{diam}K}{m_{\mathrm{far}}}
  \right\},
  \qquad \theta=1.
\]
All quantities here are finite-dimensional and come from the
Lyapunov--Schmidt map itself.

\paragraph{Singular charts.}
At a singular zero \(a_*\in Z_0\), the inverse theorem no longer gives
\(\theta=1\).  The effective replacement is a finite-jet non-flatness
estimate.  One needs a radius \(r_*>0\), an integer \(q\geq1\), and a
constant \(\kappa_*>0\) such that on the ball \(B(a_*,r_*)\) the ideal
generated by the components of
\(\Phi_0(a):=\Phi(a,\Lambda_0)\) has nontrivial normal order at most
\(q\), quantitatively bounded below by \(\kappa_*\).  Concretely, after
choosing analytic coordinates \(a=(s,n)\), where \(s\) is tangent to the
local zero set and \(n\) is normal, this means that some finite collection
of derivatives
\[
  \partial_n^\alpha \Phi_{0,j}(s,0),
  \qquad 1\leq|\alpha|\leq q,
\]
generates the normal ideal and has a controlled lower bound
\(\kappa_*\) on the relevant \(s\)-patch.

Together with a Cauchy bound
\[
  \sup_{B_\C(a_*,2r_*)}\norm{\Phi_0}\leq H_*
\]
on the complexified polydisc, quantitative Weierstrass preparation or
quantitative Lojasiewicz gives a bound of the form
\begin{equation}\label{eq:singular-loj-template}
  \operatorname{dist}(a,Z_0)
  \leq C(d,q,r_*,H_*,\kappa_*)\,
  \norm{\Phi_0(a)}^{1/q'}
\end{equation}
on \(B(a_*,r_*/2)\), where \(q'\) is controlled by \(d\) and \(q\).
The exact exponent is less important than effectivity: once
\(C\) and \(q'\) are explicit, Proposition~\ref{prop:one-chart-gap-modulus}
turns them into a gap modulus.  The true obstruction is therefore not
Fredholmness, but proving a quantitative finite-jet non-flatness bound
for the reduced analytic map at the singular charts.

\paragraph{What must be proved for the present problem.}
For the restriction problem near \(\Lambda_0\), the cleanest outcome is:
after quotienting symmetries and imposing positivity, every
level-\(\Lambda_0\) chart is regular and
\[
  Z_0/G=\{[f_0]\}.
\]
Then every chart has \(\theta=1\), and the spectral gap is controlled by
ordinary inverse-function constants.  If singular charts remain, one must
compute or estimate:
\[
  d=\dim K_*,
  \qquad q,
  \qquad r_*,
  \qquad H_*,
  \qquad \kappa_*,
  \qquad L_\lambda,
  \qquad L_a,\ L_\lambda'.
\]
These constants are produced entirely by the Fredholm reduction:
\(d\) and \(B_*\) come from the kernel and the inverse on the complement;
\(H_*\), \(M_*\), \(L_\lambda\), \(L_a\), and \(L_\lambda'\) come from
Cauchy estimates for the analytic map \(F\) and the quantitative inverse
function theorem; and \(q,\kappa_*\) are finite-jet data of the reduced
map \(\Phi_*\).  This is the finite-dimensional chart-constant problem
which replaces any tail-based strategy.

\paragraph{Dictionary of constants and expected difficulty.}
We record the precise role of the constants, since the effectiveness of
the spectral gap depends on which of them can be obtained directly from
the Fredholm reduction and which require genuinely new input.

The constant
\[
  B_*=
  \norm{\bigl(D_fF(f_*,\Lambda_0)|_{X_*}\bigr)^{-1}}
\]
is the inverse bound on the Fredholm complement after the splittings
\[
  X=K_*\oplus X_*,
  \qquad
  Y=C_*\oplus Y_*.
\]
It measures the cost of solving the infinite-dimensional equation
orthogonally to the kernel.  The integer
\[
  d=\dim K_*
\]
is the number of variables in the reduced equation
\(\Phi(a,\lambda)=0\).  These two constants are the most intrinsic
Fredholm constants.  At the constant solution they should be among the
most accessible ones, because the linearized operator is expected to
decompose into the natural symmetry and angular modes.

The constants \(H_*,r_*,M_\Phi,L_\lambda,L_a,L_\lambda'\) are analytic
chart constants.  A typical choice of \(H_*\) is a Cauchy bound on a
complexified chart,
\[
  \sup_{B_\C(a_*,2r_*)}\norm{\Phi(a,\Lambda_0)}\leq H_*.
\]
The radius \(r_*\) is the radius on which the implicit map
\(\Psi(a,\lambda)\) and the reduced map \(\Phi(a,\lambda)\) are defined
with the desired estimates.  The constant \(M_\Phi\) controls the
variation of the derivative of the reduced map,
\[
  \norm{D_a\Phi(a,\Lambda_0)-D_a\Phi(b,\Lambda_0)}
  \leq M_\Phi\norm{a-b}.
\]
The parameter Lipschitz constant \(L_\lambda\) is defined by
\[
  \norm{\Phi(a,\lambda)-\Phi(a,\Lambda_0)}
  \leq L_\lambda|\lambda-\Lambda_0|,
\]
whereas \(L_a\) and \(L_\lambda'\) control the lift from reduced
coordinates to the function space,
\[
  \norm{f(a,\Lambda_0)-f(b,\Lambda_0)}_X
  \leq L_a\norm{a-b},
  \qquad
  \norm{f(a,\lambda)-f(a,\Lambda_0)}_X
  \leq L_\lambda'|\lambda-\Lambda_0|.
\]
These constants should be obtainable, at least in principle, from
Banach-algebra estimates for the analytic norm, convolution bounds for
the fixed kernel, Cauchy estimates, and the quantitative implicit
function theorem.  They may be numerically unattractive, but they do not
encode the main geometric obstruction.

The next constants express regularity of the level-\(\Lambda_0\) zero
set.  The number
\[
  s_*=
  \inf_{a\in K\cap Z_0}
  \sigma_{\min}
  \bigl(D_a\Phi(a,\Lambda_0)|_{N_aZ_0}\bigr)
\]
is the smallest normal singular value of the reduced equation along
\(Z_0\).  If \(s_*>0\), the chart is regular in the normal directions and
one has the Lojasiewicz exponent \(\theta=1\).  The radius \(r_0\) is the
local exclusion radius around the constant orbit: within this radius,
any level-\(\Lambda_0\) solution must be the constant solution, modulo
the imposed symmetries.  Both \(s_*\) and \(r_0\) should be accessible if
one can prove that the level-\(\Lambda_0\) solution set is locally just
the constant orbit after quotienting symmetries and imposing positivity.
They are harder than the pure chart constants, but they are still
ordinary nondegeneracy constants rather than singular analytic-geometry
constants.

The residual floor
\[
  m_{\mathrm{far}}
  =
  \inf_{\substack{a\in K\\
  \operatorname{dist}(a,Z_0)\geq s_*/M_\Phi}}
  \norm{\Phi(a,\Lambda_0)}
\]
patches the local inverse estimate near \(Z_0\) into a global estimate on
the compact chart.  Although this is finite-dimensional, it is not a
cheap constant: a quantitative lower bound for \(m_{\mathrm{far}}\) is
already a finite-dimensional isolation statement away from the known
zero set.

Finally, the genuinely difficult constants are the singular
Lojasiewicz data.  The inequality
\[
  \operatorname{dist}(a,Z_0)
  \leq C_{\mathrm{Loj}}\norm{\Phi(a,\Lambda_0)}^\theta
\]
is immediate from compact real-analytic geometry, but making
\(C_{\mathrm{Loj}}\) and \(\theta\) explicit is easy only in regular
charts.  In singular charts one needs a finite order \(q\) and a
non-flatness lower bound \(\kappa_*\), schematically
\[
  \max_{1\leq|\alpha|\leq q,j}
  |\partial_n^\alpha\Phi_{0,j}(s,0)|
  \geq \kappa_*,
\]
on the relevant tangential patch.  Together with \(H_*\) and \(r_*\),
these data feed quantitative Weierstrass preparation or quantitative
Lojasiewicz estimates and produce an exponent of the form
\(\theta=1/q'\).  This is the hardest part of the program if singular
charts cannot be ruled out.

Thus the practical hierarchy is as follows.  The constants
\[
  d,\ B_*,\ H_*,\ r_*,\ M_\Phi,\ L_\lambda,\ L_a,\ L_\lambda'
\]
are bookkeeping constants of the Fredholm reduction and analytic
implicit function theorem.  The constants
\[
  s_*,\ r_0,\ m_{\mathrm{far}}
\]
are quantitative isolation and nondegeneracy constants for the
level-\(\Lambda_0\) solution set.  The constants
\[
  q,\ \kappa_*,\ C_{\mathrm{Loj}},\ \theta
\]
are the singular analytic-geometry constants.  The preferred route is
therefore to prove that, after quotienting symmetries and imposing
positivity, the relevant level-\(\Lambda_0\) charts are regular.  Then
\(\theta=1\), \(C_{\mathrm{Loj}}\) is controlled by ordinary inverse
function constants, and the gap follows from
\[
  L_a C_{\mathrm{Loj}}L_\lambda^\theta\eta_*^\theta
  +L_\lambda'\eta_*<r_0.
\]
If singular charts remain, the program is still finite-dimensional, but
one must supply effective finite-jet non-flatness estimates for
\(\Phi(\cdot,\Lambda_0)\).

\subsection{Residual functional and distance to zero sets}
\label{sec:residual-functional}

Another way to package the preceding plan is to study a residual
functional.  Let $X$ be a normed target space in which the residual is
defined and in which small residual has geometric meaning.  Formally,
one would like to use
\begin{equation}\label{eq:residual-functional}
  \mathcal E_X(f,\lambda)
  :=
  \norm{T(f)-\lambda f}_X^2,
  \qquad \lambda\in\R,\qquad \norm f_2=1.
\end{equation}
Zeros of $\mathcal E_X$ are precisely nonlinear eigenpairs.  If $X$ is
chosen so that $T(f)-\lambda f$ is analytic in $f$ and $\lambda$, then
the zero set is an analytic variety, at least after restricting to the
analytic regularity class supplied by the bootstrap.

There is a small functional-analytic caveat.  If $f$ is only assumed to
lie in $L^2$, then a strong analytic norm $X$ need not contain the term
$\lambda f$.  Thus one should either take $X=L^2$ or a negative Sobolev
space for the global residual, or first use the a priori bootstrap to
place all relevant solutions in a fixed analytic Banach space
$\mathcal A_\rho$ and then define $\mathcal E_X$ there.  In the present
route, finite-dimensionality enters only locally, through Fredholm
Lyapunov--Schmidt reduction, not through a global Fourier cut-off.

The essential quantitative question is then:
\[
  \text{how can one bound the distance to } Z(P):=\{P=0\}
  \text{ directly from } \norm{P}?
\]
The standard answer is a quantitative Lojasiewicz inequality.  In a
Fredholm chart, the infinite-dimensional equation is reduced to
\[
  \Phi(a,\lambda)=0,
  \qquad a\in U\subset\R^d,
\]
where $\Phi$ is real-analytic.  On a compact set $K\Subset U$, one seeks
explicit constants $C,\theta>0$ such that
\begin{equation}\label{eq:lojasiewicz-distance}
  \operatorname{dist}(a,Z(\Phi_{\Lambda_0})\cap K)
  \leq C\norm{\Phi(a,\Lambda_0)}^\theta,
  \qquad a\in K,
\end{equation}
where $\Phi_{\Lambda_0}(a):=\Phi(a,\Lambda_0)$.  If
$a\in Z(\Phi_\lambda)$, then
\[
  \Phi(a,\Lambda_0)
  =
  \Phi(a,\Lambda_0)-\Phi(a,\lambda),
\]
so an explicit bound on $\partial_\lambda\Phi$ gives
\[
  \norm{\Phi(a,\Lambda_0)}
  \leq L_\lambda|\lambda-\Lambda_0|.
  \]
Substituting in \eqref{eq:lojasiewicz-distance} yields
\begin{equation}\label{eq:zero-set-distance-from-loj}
  \operatorname{dist}\bigl(a,Z(\Phi_{\Lambda_0})\cap K\bigr)
  \leq C L_\lambda^\theta |\lambda-\Lambda_0|^\theta.
\end{equation}
This is exactly the local form of the effective upper-semicontinuity
modulus \eqref{eq:effective-upper-semicontinuity}.

Thus the residual-functional route reduces the nonlinear spectral gap to
three explicit analytic estimates:
\begin{enumerate}[leftmargin=2em]
  \item a bootstrap placing all relevant eigenfunctions in a compact
  analytic class, together with finitely many gauge-fixed Fredholm charts;
  \item in each Lyapunov--Schmidt chart, an effective Lojasiewicz or
  Weierstrass-preparation bound for the analytic residual
  $\Phi(\cdot,\Lambda_0)$ on the relevant compact set;
  \item effective local rigidity of the zero set at level $\Lambda_0$,
  ideally $\mathcal S_{\Lambda_0}/G=\{[f_0]\}$ in the positive gauge.
\end{enumerate}
If these are available, then the spectral gap is quantified by choosing
$|\lambda-\Lambda_0|$ so small that the Lojasiewicz distance to
$\mathcal S_{\Lambda_0}$ is below the local isolation radius~$r_0$.

\subsection{A projected-equation route to the gap}
\label{sec:projected-gap-route}

We now record a more analytic route which avoids a blind numerical search
over the annulus.  The guiding principle is that a critical point with
\(\Lambda\) close to \(\Lambda_0\) should be forced into the local
nondegeneracy ball around~\(f_0\).  Once this is known, the
Newton--Kantorovich/local uniqueness argument excludes it unless it is
exactly~\(f_0\).

Throughout this subsection we work in the real Hilbert space underlying
\(L^2(S^1)\), with inner product \(\re\ip{\cdot}{\cdot}\), and write
\[
  T(f):=E^*(|Ef|^4Ef).
\]
Thus critical points satisfy
\[
  T(f)=\Lambda f,\qquad \norm f_2=1,
  \qquad
  \Lambda=J(f):=\norm{Ef}_6^6.
\]

\paragraph{Gauge and coercive slice.}
Let \(f_0=(2\pi)^{-1/2}\).  The tangent space to the symmetry orbit at
\(f_0\) is generated by the phase direction \(if_0\) and the two
translation directions, which are the \(n=\pm1\) Fourier modes.  For the
positive problem the phase is fixed, and we assume that the translation
gauge has been chosen so that the perturbation is orthogonal to the
translation directions.  In this gauge write
\begin{equation}\label{eq:af0-xi-decomp}
  f=af_0+\xi,
  \qquad
  a=\re\ip{f}{f_0}>0,
  \qquad
  \xi\perp f_0,\quad \xi\perp T_{f_0}(G\cdot f_0).
\end{equation}
Since \(\norm f_2=1\), we have
\[
  a^2+\sigma^2=1,\qquad \sigma:=\norm{\xi}_2.
\]
The point of the gauge is that \(\xi\) belongs to the spectral subspace
on which the linearised operator has a genuine gap.

\paragraph{Linear spectral input.}
Let \(K:=DT(f_0)\).  On the coercive slice define
\[
  \mu_0
  :=
  \sup_{\substack{\xi\perp f_0,\ \xi\perp T_{f_0}(G\cdot f_0)\\
                  \norm{\xi}_2=1}}
  \re\ip{K\xi}{\xi}.
\]
By the Fourier decomposition in Proposition~\ref{prop:eigenvalues},
\[
  \mu_0=\max_{n\geq2} 5\mathcal I_n.
\]
The certified Bessel gap gives, for example,
\[
  \mu_0\leq 0.6\Lambda_0.
\]
For the sharp local constant one expects
\[
  \mu_0=5\mathcal I_2,
  \qquad
  \Lambda_0-\mu_0\approx0.452\Lambda_0.
\]
Only the upper bound \(\mu_0<\Lambda_0\) is needed for the abstract
argument below.

\paragraph{Finite Taylor expansion of \(T\).}
Since \(T\) is a homogeneous polynomial map of degree five on the real
Hilbert space, the Taylor expansion around~\(f_0\) terminates:
\begin{equation}\label{eq:T-taylor-project}
  T(f_0+h)
  =
  \Lambda_0f_0
  +Kh
  +\sum_{j=2}^{5}\frac1{j!}D^jT(f_0)[h^j].
\end{equation}
For \(f=af_0+\xi\) it is slightly cleaner to factor out the exact
homogeneity:
\[
  T(af_0+\xi)
  =
  a^5T\left(f_0+\frac{\xi}{a}\right)
  \qquad(a>0).
\]
Using \eqref{eq:T-taylor-project}, this gives
\begin{equation}\label{eq:T-af0-xi}
  T(af_0+\xi)
  =
  a^5\Lambda_0f_0
  +a^4K\xi
  +\sum_{j=2}^{5}a^{5-j}\frac1{j!}D^jT(f_0)[\xi^j].
\end{equation}

\paragraph{Remainder constants.}
For \(2\leq j\leq5\), define the coercive-slice multilinear constants
\begin{equation}\label{eq:Cj-def}
  C_j
  :=
  \sup_{\substack{\xi\perp f_0,\ \xi\perp T_{f_0}(G\cdot f_0)\\
                  \norm{\xi}_2=1}}
  \left|
    \re\ip{\frac1{j!}D^jT(f_0)[\xi^j]}{\xi}
  \right|.
\end{equation}
More generally, one may replace \(C_j\) by the operator norm of the
\((j+1)\)-linear form
\[
  (\xi_1,\ldots,\xi_j,\zeta)
  \mapsto
  \re\ip{\frac1{j!}D^jT(f_0)[\xi_1,\ldots,\xi_j]}{\zeta}
\]
on the coercive slice.  The diagonal constants in \eqref{eq:Cj-def} are
the smallest constants needed for the following scalar proof.

Each \(C_j\) is in principle an explicit finite/infinite sum of Bessel
product integrals.  For example, \(C_2\) involves integrals of the form
\[
  \int_0^\infty J_0(r)^3J_{n_1}(r)J_{n_2}(r)J_{n_3}(r)\,r\,dr,
\]
with angular selection rules.  The higher \(C_j\)'s involve the same
sixfold Bessel product family appearing in the linear and multilinear
expansions at the constant.  The analytic proof only needs upper bounds
for these constants; any computer assistance would enter only as
certified evaluation of explicit Bessel-product norms.

\paragraph{The projected equation.}
Assume now that \(f=af_0+\xi\) is a positive normalised critical point in
the above gauge:
\[
  T(f)=\Lambda f.
\]
Take the real inner product with~\(\xi\).  Since \(\xi\perp f_0\), the
constant term in \eqref{eq:T-af0-xi} drops out and we obtain
\begin{align}
  \Lambda\sigma^2
  &=
  a^4\re\ip{K\xi}{\xi}
  +
  \sum_{j=2}^{5}a^{5-j}
  \re\ip{\frac1{j!}D^jT(f_0)[\xi^j]}{\xi}
  \notag\\
  &\leq
  a^4\mu_0\sigma^2
  +
  C_2a^3\sigma^3
  +
  C_3a^2\sigma^4
  +
  C_4a\sigma^5
  +
  C_5\sigma^6.
  \label{eq:projected-master-ineq}
\end{align}
This is the central inequality.  It is purely analytic and contains only
the linear spectral constant \(\mu_0\) and the explicit multilinear
constants \(C_2,\ldots,C_5\).

\begin{proposition}[Projected-equation exclusion criterion]
  \label{prop:projected-exclusion}
  Let \(0<r_0<1\) be a local uniqueness radius around \(f_0\) in the
  positive gauge, and let \(\eta>0\).  Suppose that for every
  \(\sigma\in[r_0,1]\), with \(a=\sqrt{1-\sigma^2}\), one has
  \begin{equation}\label{eq:projected-exclusion-criterion}
    (\Lambda_0-\eta-a^4\mu_0)\sigma^2
    >
    C_2a^3\sigma^3
    +C_3a^2\sigma^4
    +C_4a\sigma^5
    +C_5\sigma^6 .
  \end{equation}
  Then there is no positive normalised critical point \(f\neq f_0\),
  modulo symmetries, with
  \[
    |\Lambda-\Lambda_0|\leq\eta.
  \]
  In particular, the nonlinear critical-value gap around \(\Lambda_0\)
  satisfies
  \[
    \gamma_0\geq\eta.
  \]
\end{proposition}

\begin{proof}
  Suppose \(T(f)=\Lambda f\), \(\norm f_2=1\), \(f\geq0\), and
  \(|\Lambda-\Lambda_0|\leq\eta\).  Put \(f\) in the positive translation
  gauge and decompose it as in \eqref{eq:af0-xi-decomp}.  If
  \(\sigma<r_0\), local uniqueness gives \(f=f_0\), hence
  \(\Lambda=\Lambda_0\).  Otherwise \(\sigma\in[r_0,1]\).

  Since \(\Lambda\geq\Lambda_0-\eta\), the left side of
  \eqref{eq:projected-master-ineq} is bounded below by
  \((\Lambda_0-\eta)\sigma^2\).  Hence
  \[
    (\Lambda_0-\eta)\sigma^2
    \leq
    a^4\mu_0\sigma^2
    +C_2a^3\sigma^3
    +C_3a^2\sigma^4
    +C_4a\sigma^5
    +C_5\sigma^6,
  \]
  contradicting \eqref{eq:projected-exclusion-criterion}.  Thus no such
  \(f\) exists outside the local ball, and the local ball contains only
  \(f_0\).
\end{proof}

\begin{remark}[Why this proves exact-level uniqueness]
  Taking \(\eta=0\) in \eqref{eq:projected-exclusion-criterion} gives a
  direct uniqueness criterion at the exact value~\(\Lambda_0\):
  \[
    T(f)=\Lambda_0f,\quad \norm f_2=1,\quad f\geq0
    \quad\Longrightarrow\quad
    f=f_0
  \]
  modulo symmetries.  Taking \(\eta>0\) gives the stronger quantitative
  gap statement.
\end{remark}

\begin{remark}[A usable sufficient form]
  Since \(a\leq1\), a cruder but simpler sufficient condition is
  \begin{equation}\label{eq:projected-crude}
    \Lambda_0-\eta-\mu_0
    >
    C_2\sigma+C_3\sigma^2+C_4\sigma^3+C_5\sigma^4
    \qquad(r_0\leq\sigma\leq1).
  \end{equation}
  This form is usually too wasteful for large~\(\sigma\), but it shows the
  mechanism clearly: the linear spectral gap \(\Lambda_0-\mu_0\) must
  dominate the nonlinear remainders on the complement of the local ball.
  A sharper implementation should use
  \eqref{eq:projected-exclusion-criterion}, because the factor
  \(a=\sqrt{1-\sigma^2}\) suppresses the lower-order nonlinear terms away
  from the constant mode.
\end{remark}

\paragraph{What remains to make the criterion effective.}
The proof of Proposition~\ref{prop:projected-exclusion} reduces the
nonlinear gap problem to the following finite list of quantitative inputs:
\begin{enumerate}[leftmargin=2em]
  \item a certified local uniqueness radius \(r_0\), obtained from
    \(\nu_0\) and the Newton--Kantorovich Lipschitz constant;
  \item the linear coercive constant
    \[
      \mu_0=\sup_{\xi\in\mathcal S,\ \norm{\xi}=1}\re\ip{K\xi}{\xi},
    \]
    where \(\mathcal S\) is the slice orthogonal to \(f_0\) and the
    symmetry directions;
  \item certified upper bounds for the multilinear constants
    \(C_2,\ldots,C_5\);
  \item a one-variable verification of
    \eqref{eq:projected-exclusion-criterion} on
    \(\sigma\in[r_0,1]\).
\end{enumerate}
This is the desired ``straight proof'' architecture: the argument itself is
analytic, while numerical work is confined to explicit Bessel-product
constants and a final one-dimensional inequality.


\section{Options for Completing the Sharp Constant}

Two paths are available, depending on the desired strength of the result.

\paragraph{Path 1: Identify $\mathcal R$ (the sharp constant).}
\begin{enumerate}[leftmargin=2em]
  \item \textbf{Verify Bessel integrals.}  Check
    $\mathcal J_n\neq\mathcal J_0$ and $5\mathcal J_n\neq\mathcal J_0$
    for $n=2,3,\ldots,n_*$ by interval arithmetic.  This establishes
    nondegeneracy at the constant solution.  \emph{Routine numerical
    computation.}

  \item \textbf{Restrict the search to nonnegative eigenfunctions.}
    By Lemma~\ref{lem:positive-reduction}, any extremiser can be replaced
    by a nonnegative extremiser, possibly in a different shape.  Therefore,
    for the sharp constant, one only needs to retain those candidate
    values~$\Lambda_j$ for which the nonlinear eigenvalue equation has at
    least one nonnegative solution.  The positive gauge also removes exact
    translation ambiguity by Lemma~\ref{lem:positive-gauge}, and
    Theorem~\ref{thm:positive-values-finite} proves that the retained
    critical values in any fixed positive window are finite.

  \item \textbf{Show the constant is the unique maximiser.}
    Theorem~\ref{thm:maximiser} gives existence of a maximiser.
    If one can prove that the constant function $f_0=(2\pi)^{-1/2}$
    is the unique maximiser modulo~$G$---e.g., by a rearrangement
    or symmetrisation argument, or by ruling out non-radial
    maximisers via the eigenvalue analysis---then
    $\mathcal R=\norm{Ef_0}_{L^6}$ and the programme is complete.
    \emph{This bypasses compactness of the full critical set.}
\end{enumerate}

\paragraph{Path 2: Finiteness of all critical values above
$\varepsilon$ (the full structural result).}
\begin{enumerate}[leftmargin=2em]
  \item Bessel integrals (same as above).

  \item \textbf{Rule out multi-profile splitting
    (Theorem~\ref{thm:compactness-mod-G}).}
    The profile decomposition is in place
    (Theorem~\ref{thm:profile-decomp}), and each profile inherits the
    Euler--Lagrange equation
    (Proposition~\ref{prop:profile-critical}).  The open step is to
    show that no critical sequence admits more than one profile.
    Section~\ref{sec:compactness} identifies three general strategies:
    descending induction, Pohozaev--Morawetz monotonicity, and unique
    continuation for cross terms.  \emph{Requires new analytic input.}

  \item \textbf{Enumerate non-constant shapes (if any).}  Either prove
    that the constant is the only critical shape, or verify
    nondegeneracy at all other shapes.
\end{enumerate}

% =======================================================================
% PART III: EXTENSIONS AND SUPPLEMENTARY MATERIAL
% =======================================================================



\end{document}
